\documentclass[
  journal=small,
  manuscript=research-article,
  year=2026
]{cup-journal}

\usepackage{booktabs}
\usepackage{longtable}
\usepackage{array}
\usepackage{amsmath}
\usepackage{url}
\usepackage{xcolor}
\usepackage[hidelinks]{hyperref}

\newcolumntype{L}[1]{>{\raggedright\arraybackslash}p{#1}}
\sloppy

\title{Supplementary Appendix: Constitutional-Review Design-Bundle Stress Testing}
\author{Anonymous Author}
\affiliation{Manuscript prepared for anonymous review}
\date{}

\begin{document}
\maketitle

\section*{Supplementary Appendix Scope}

This appendix preserves technical detail that is too large for the main Journal of Law and Courts manuscript: fixed model weights, full main-comparison results, sensitivity results, calibration source-observation status, empirical data-development files, and aggregate artifact inventories. It is organized around the manuscript's evidentiary split. Stylized archetype presets are institutional mappings, mechanism scenarios are counterfactual design tests, and candidate source rows enter source-range checks only after primary-source verification. The main manuscript presents the project as a comparative stress-testing framework centered on veto-relocation risk; emergency-process reforms, appointment changes, councils, overrides, weak-form review, and front-end review mechanisms are compared under the same generated worlds. The replication archive includes source code, tests, configuration inputs, aggregate outputs, replication manifests, and build commands. Full case-level and interval outputs are generated by the same commands and should be deposited with the final public replication archive or supplied separately if requested during review.

\section*{Historical Benchmark and Robustness v1}

This study is additional to the original campaigns below. Its frozen protocol, exposure audit, first evaluation, analysis contract, and source hashes are retained in the replication archive. The empirical test concerns the composition of 66 citation-centered SCDB disputes in the 2025 term, not petitions or constitutional merits cases alone. SCDB 2025 Release 01 remains the training source; 2026 Release 01 is used for the new term and a separate full-field revision audit. The audit finds two additions and ten changed matched historical records, with 9,331 matched records unchanged. Older-period revisions never overwrite the training source. Other remains in the category denominator, and the legacy calibration predicates remain unchanged.

The protocol preceded acquisition of the new-term case records, and the four forecast vectors were locked before that acquisition. Development checks (2005--2014) and backtests (2015--2024) are explicitly retrospective because old aggregate calibration had already exposed those periods. This is a scoped new-data test, not blindness to public events. The primary recency-minus-fixed Brier comparison shows no clear conditional advantage; all secondary comparisons remain visible in Table~\ref{tab:supp-historical-scores}.

% Auto-generated by paper/scripts/generate_historical_tables.py
{\small
\setlength{\tabcolsep}{3pt}
\begin{longtable}{@{}L{0.55\linewidth}rr@{}}
\caption{Observed 2025 test categories, without denominator exclusions.}\label{tab:supp-historical-categories}\\
\toprule
Category & Disputes & Share \\
\midrule
\endfirsthead
\toprule
Category & Disputes & Share \\
\midrule
\endhead
Speech & 2 & 0.0303 \\
Equality & 6 & 0.0909 \\
Criminal Procedure & 20 & 0.3030 \\
Federalism & 3 & 0.0455 \\
Election Law & 4 & 0.0606 \\
Emergency Powers & 2 & 0.0303 \\
Administrative State & 16 & 0.2424 \\
Other & 13 & 0.1970 \\
\bottomrule
\end{longtable}
\par\noindent\begin{minipage}{0.98\linewidth}\footnotesize The eight categories form a mutually exclusive partition of all 66 citation-centered disputes. Other remains in the denominator. Observed test shares were not used as simulation profiles; the four simulation profiles are the training-only forecasts retained in the replication inputs.\end{minipage}
}

% Auto-generated by paper/scripts/generate_historical_tables.py
{\small
\setlength{\tabcolsep}{3pt}
\begin{longtable}{@{}L{0.33\linewidth}rrrr@{}}
\caption{Locked training-only forecast probabilities.}\label{tab:supp-historical-forecasts}\\
\toprule
Category & Fixed & Recency & Expanding & Previous \\
\midrule
\endfirsthead
\toprule
Category & Fixed & Recency & Expanding & Previous \\
\midrule
\endhead
Speech & 0.0700 & 0.0592 & 0.0726 & 0.0662 \\
Equality & 0.1580 & 0.1636 & 0.1521 & 0.1250 \\
Criminal Procedure & 0.2360 & 0.2189 & 0.2251 & 0.1838 \\
Federalism & 0.0565 & 0.0449 & 0.0559 & 0.0074 \\
Election Law & 0.0262 & 0.0266 & 0.0264 & 0.0074 \\
Emergency Powers & 0.0060 & 0.0099 & 0.0011 & 0.0074 \\
Administrative State & 0.2190 & 0.2506 & 0.2171 & 0.2868 \\
Other & 0.2283 & 0.2263 & 0.2497 & 0.3162 \\
\bottomrule
\end{longtable}
\par\noindent\begin{minipage}{0.98\linewidth}\footnotesize Probabilities are displayed rounded; full precision is retained in the frozen JSON. Recency uses 2010--2024 (1,059 disputes); expanding uses 1946--2024 (9,341); previous uses 2024 (64). The fixed vector was already informed by legacy aggregate calibration. The test categories in the preceding table are never substituted into these profiles.\end{minipage}
}

% Auto-generated by paper/scripts/generate_historical_tables.py
{\small
\setlength{\tabcolsep}{3pt}
\begin{longtable}{@{}L{0.25\linewidth}L{0.21\linewidth}rL{0.28\linewidth}@{}}
\caption{All frozen empirical paired-score comparisons.}\label{tab:supp-historical-scores}\\
\toprule
Comparison & Score & Difference & 95\% interval \\
\midrule
\endfirsthead
\toprule
Comparison & Score & Difference & 95\% interval \\
\midrule
\endhead
Recency - fixed & multiclass-brier & 0.001704 & [-0.007500, 0.010336] \\
Recency - fixed & log-loss & -0.011810 & [-0.044064, 0.018033] \\
Recency - fixed & total-variation & -0.010039 & [-0.032769, 0.026862] \\
Recency - expanding & multiclass-brier & -0.001465 & [-0.011552, 0.008137] \\
Recency - expanding & log-loss & -0.063817 & [-0.167137, 0.015321] \\
Recency - expanding & total-variation & -0.027576 & [-0.036322, 0.031648] \\
Recency - previous & multiclass-brier & -0.021949 & [-0.045585, 0.001273] \\
Recency - previous & log-loss & -0.144563 & [-0.276378, -0.031456] \\
Recency - previous & total-variation & -0.094387 & [-0.118072, -0.003049] \\
\bottomrule
\end{longtable}
\par\noindent\begin{minipage}{0.98\linewidth}\footnotesize All scores are lower-is-better. The primary endpoint is recency-minus-fixed Brier score. Secondary advantages against the previous-term baseline do not replace that endpoint. Paired-case bootstrap uncertainty holds the forecasts fixed within the 66-dispute term.\end{minipage}
}


Training uncertainty is separately reported by resampling the fifteen training terms with fixed age weights. Retrospective results use equal-target-term means and paired target-term cluster intervals; case-weighted means are also retained in the CSV. These quantities are not the new-term paired-case uncertainty or the institutional run-bootstrap uncertainty. The empirical report includes every development/backtest model and score, category mapping statuses, and the historical revision audit. No favorable retrospective or secondary result replaces the frozen primary comparison.

\subsection*{Institutional grid and chronology}

Four locked forecast profiles cross the thirteen contexts in Table~\ref{tab:supp-historical-contexts} and eight designs. Each cell has 120 runs, 80 cases per run, and four review periods. Within a run, all designs receive the same complete case/object docket. Across forecast/context cells, matching run numbers do not imply identical dockets or independent real-world environments. Institutional voting streams are design-specific; response thinning uses a separate recorded stream. The response interventions scale credibility before downstream evaluation and independently thin positive response events. A factor of one reproduces legacy behavior exactly; half response need not halve downstream outcomes.

The procedure intervention changes reasoned review to fast emergency handling without changing its other configuration fields. The structure intervention changes council structure to full-court review while retaining its other intake and archetype terms. Half and off variants retain every non-response weak-form term. These restricted interventions expose model dependence; a null upstream effect is not corrected by silently broadening the intervention.

% Auto-generated by paper/scripts/generate_historical_tables.py
{\small
\setlength{\tabcolsep}{3pt}
\begin{longtable}{@{}L{0.55\linewidth}rr@{}}
\caption{One-at-a-time context grid.}\label{tab:supp-historical-contexts}\\
\toprule
World input & Low & High \\
\midrule
\endfirsthead
\toprule
World input & Low & High \\
\midrule
\endhead
emergencyPressure & 0.1 & 0.6 \\
rightsThreatRate & 0.2 & 0.6 \\
appointmentPolarization & 0.25 & 0.85 \\
legislativeConflict & 0.15 & 0.75 \\
implementationCapacity & 0.25 & 0.85 \\
partyFragmentation & 0.2 & 0.8 \\
\bottomrule
\end{longtable}
\par\noindent\begin{minipage}{0.98\linewidth}\footnotesize Baseline plus twelve one-at-a-time settings gives thirteen contexts per forecast. Other world inputs remain at their baseline values. This finite grid assigns no probability to real political environments and is not a full interaction design.\end{minipage}
}


The original protocol fixed the grid, component families, intervention factors, 2,000-replicate resampling scheme, 0.01 practical threshold, and normative exercise. The signed measurement contract was committed after simulation and raw audit, before comparative estimates. It is an implementation interpretation of the original protocol, not an independent preregistration. Three additional diagnostics are marked separately. The first calculation later exposed a low-emergency conditional-rate reversal; pooled event diagnostics were added after that estimate without changing any original contrast or leadership result.

\subsection*{Complete claims, interventions, and uncertainty}

The base contrasts are reasoned minus current, council minus current, and weak form minus current. Each point estimate averages 120 within-run differences. Resampling whole paired runs preserves within-run dependence; the same sampled-index matrix is reused across all cells and metrics. Intervals are conditional Monte Carlo intervals, not model-specification intervals. The full analysis contains 3,692 contrasts across 52 forecast/context cells. Tables~\ref{tab:supp-historical-claims} and \ref{tab:supp-historical-ablations} include all components and interventions without filtering on their signs. Table~\ref{tab:supp-historical-intervals} shows all 71 reference-cell intervals; the complete CSV contains every other cell and its standard errors.

% Auto-generated by paper/scripts/generate_historical_tables.py
{\small
\setlength{\tabcolsep}{3pt}
\begin{longtable}{@{}L{0.12\linewidth}L{0.25\linewidth}L{0.21\linewidth}L{0.19\linewidth}r@{}}
\caption{Complete component-level claim matrix.}\label{tab:supp-historical-claims}\\
\toprule
Claim & Component & Difference range & Pattern & CI +/$-$ \\
\midrule
\endfirsthead
\toprule
Claim & Component & Difference range & Pattern & CI +/$-$ \\
\midrule
\endhead
Emergency & Reasons per order & -0.0255 to 0.8798 & sign-reversing & 48/1 \\
Emergency & Merits follow-up & 0.0083 to 1.0000 & mixed/weak & 49/0 \\
Emergency & Review-time input & 0.0000 to 0.0000 & mixed/weak & 0/0 \\
Emergency & Total cost index & 0.0095 to 0.0137 & mixed/weak & 52/0 \\
Emergency & Emergency irregularity & -0.8280 to -0.0456 & stable & 0/52 \\
Emergency & Delay-cost index$^\dagger$ & 0.0029 to 0.0085 & descriptive-only & 52/0 \\
Screening & Direct court cost & 0.1358 to 0.1375 & descriptive-only & 52/0 \\
Screening & Upstream cost & 0.0004 to 0.0022 & mixed/weak & 52/0 \\
Screening & Capacity strain & 0.0631 to 0.0793 & descriptive-only & 52/0 \\
Screening & Total cost index & 0.0828 to 0.0860 & descriptive-only & 52/0 \\
Screening & Intake acceptance & -0.0008 to -0.0002 & descriptive-only & 0/52 \\
Screening & Council screening & 1.0000 to 1.0000 & stable & 52/0 \\
Screening & Administrative load$^\dagger$ & 0.2462 to 0.2683 & descriptive-only & 52/0 \\
Weak form & Legislative response & 0.1496 to 0.5448 & stable & 52/0 \\
Weak form & Response credibility & 0.0443 to 0.0905 & stable & 52/0 \\
Weak form & Timely replies & 0.0000 to 0.3279 & mixed/weak & 40/0 \\
Weak form & Rights index & -0.0055 to 0.0598 & mixed/weak & 48/4 \\
Weak form & Compliance events & 0.1923 to 0.3909 & stable & 52/0 \\
Weak form & Responsiveness & 0.0467 to 0.1088 & stable & 52/0 \\
Weak form & Veto-risk index & -0.3813 to -0.2907 & stable & 0/52 \\
Weak form & Total cost index & -0.0947 to -0.0844 & descriptive-only & 0/52 \\
Weak form & Reply delay$^\dagger$ & 0.3607 to 0.4895 & descriptive-only & 52/0 \\
\bottomrule
\end{longtable}
\par\noindent\begin{minipage}{0.98\linewidth}\footnotesize All 22 components are shown. $\dagger$ Additional diagnostics, not replacement endpoints. Stable requires every signed base estimate beyond the expected strict 0.01 threshold; sign-reversing requires both practical signs. Unsigned rows are descriptive-only. CI counts are unadjusted intervals above/below zero, out of 52. Duration is a shared input; conditional reasons, follow-up and reply timing depend on eligible-event populations. Mechanism activation and constructed indices are not independent validation.\end{minipage}
}

% Auto-generated by paper/scripts/generate_historical_tables.py
{\small
\setlength{\tabcolsep}{3pt}
\begin{longtable}{@{}L{0.26\linewidth}L{0.27\linewidth}L{0.24\linewidth}r@{}}
\caption{All mechanism and persistence contrasts.}\label{tab:supp-historical-ablations}\\
\toprule
Contrast & Component & Difference range & CI +/$-$ \\
\midrule
\endfirsthead
\toprule
Contrast & Component & Difference range & CI +/$-$ \\
\midrule
\endhead
Fast - reasoned & Reasons per order & -0.7348 to 0.1774 & 4/48 \\
Fast - reasoned & Merits follow-up & -1.0000 to -0.0083 & 0/49 \\
Fast - reasoned & Review-time input & 0.0000 to 0.0000 & 0/0 \\
Fast - reasoned & Total cost index & -0.0083 to -0.0042 & 0/52 \\
Fast - reasoned & Emergency irregularity & 0.0458 to 0.7970 & 52/0 \\
Fast - reasoned & Delay-cost index & -0.0068 to -0.0013 & 0/52 \\
Full court - council & Direct court cost & -0.0973 to -0.0971 & 0/52 \\
Full court - council & Upstream cost & -0.0001 to 0.0001 & 1/1 \\
Full court - council & Capacity strain & -0.0438 to -0.0425 & 0/52 \\
Full court - council & Total cost index & -0.0504 to -0.0501 & 0/52 \\
Full court - council & Intake acceptance & 0.0000 to 0.0000 & 1/1 \\
Full court - council & Council screening & -1.0000 to -1.0000 & 0/52 \\
Full court - council & Administrative load & -0.2410 to -0.2391 & 0/52 \\
Half - weak form & Legislative response & -0.2687 to -0.0744 & 0/52 \\
Half - weak form & Response credibility & -0.3258 to -0.2742 & 0/52 \\
Half - weak form & Timely replies & -0.3229 to 0.0000 & 0/40 \\
Half - weak form & Rights index & -0.0191 to 0.0014 & 2/48 \\
Half - weak form & Compliance events & -0.0697 to 0.0136 & 1/17 \\
Half - weak form & Responsiveness & -0.0401 to -0.0098 & 0/52 \\
Half - weak form & Veto-risk index & 0.0121 to 0.0433 & 52/0 \\
Half - weak form & Total cost index & -0.0026 to -0.0006 & 0/52 \\
Half - weak form & Reply delay & 0.0708 to 0.0860 & 52/0 \\
Off - weak form & Legislative response & -0.5448 to -0.1496 & 0/52 \\
Off - weak form & Response credibility & -0.6521 to -0.5483 & 0/52 \\
Off - weak form & Timely replies & -0.3279 to 0.0000 & 0/40 \\
Off - weak form & Rights index & -0.0328 to 0.0036 & 4/48 \\
Off - weak form & Compliance events & -0.0952 to 0.0026 & 0/22 \\
Off - weak form & Responsiveness & -0.0807 to -0.0189 & 0/52 \\
Off - weak form & Veto-risk index & 0.0241 to 0.0872 & 52/0 \\
Off - weak form & Total cost index & -0.0051 to -0.0014 & 0/52 \\
Off - weak form & Reply delay & -0.4895 to -0.3607 & 0/52 \\
Half - current & Legislative response & 0.0752 to 0.2773 & 52/0 \\
Half - current & Response credibility & -0.2676 to -0.2125 & 0/52 \\
Half - current & Timely replies & 0.0000 to 0.0050 & 3/0 \\
Half - current & Rights index & -0.0207 to 0.0610 & 44/8 \\
Half - current & Compliance events & 0.1776 to 0.3873 & 52/0 \\
Half - current & Responsiveness & 0.0239 to 0.0974 & 52/0 \\
Half - current & Veto-risk index & -0.3383 to -0.2780 & 0/52 \\
Half - current & Total cost index & -0.0953 to -0.0852 & 0/52 \\
Half - current & Reply delay & 0.4444 to 0.5624 & 52/0 \\
Off - current & Legislative response & 0.0000 to 0.0000 & 0/0 \\
Off - current & Response credibility & -0.5923 to -0.4889 & 0/52 \\
Off - current & Timely replies & 0.0000 to 0.0000 & 0/0 \\
Off - current & Rights index & -0.0323 to 0.0634 & 44/8 \\
Off - current & Compliance events & 0.1625 to 0.3842 & 52/0 \\
Off - current & Responsiveness & -0.0002 to 0.0845 & 48/0 \\
Off - current & Veto-risk index & -0.3031 to -0.2654 & 0/52 \\
Off - current & Total cost index & -0.0964 to -0.0861 & 0/52 \\
Off - current & Reply delay & 0.0000 to 0.0000 & 0/0 \\
\bottomrule
\end{longtable}
\par\noindent\begin{minipage}{0.98\linewidth}\footnotesize Ranges span the same 52 cells. Fast changes only emergency procedure; full court changes only council review structure. Half/off scale response credibility and thin positive response events. Half/off minus current are persistence checks. No replies is not a late or instantaneous reply; reply-conditioned comparisons are not population-matched.\end{minipage}
}

% Auto-generated by paper/scripts/generate_historical_tables.py
{\small
\setlength{\tabcolsep}{3pt}
\begin{longtable}{@{}L{0.25\linewidth}L{0.25\linewidth}rL{0.26\linewidth}@{}}
\caption{Paired intervals in the fixed-generator, baseline-context cell.}\label{tab:supp-historical-intervals}\\
\toprule
Contrast & Component & Difference & 95\% MC interval \\
\midrule
\endfirsthead
\toprule
Contrast & Component & Difference & 95\% MC interval \\
\midrule
\endhead
Reasoned - current & Reasons per order & 0.8443 & [0.8147, 0.8725] \\
Reasoned - current & Merits follow-up & 0.9917 & [0.9750, 1.0000] \\
Reasoned - current & Review-time input & 0.0000 & [0.0000, 0.0000] \\
Reasoned - current & Total cost index & 0.0098 & [0.0097, 0.0100] \\
Reasoned - current & Emergency irregularity & -0.3199 & [-0.3276, -0.3120] \\
Reasoned - current & Delay-cost index & 0.0082 & [0.0080, 0.0084] \\
Council - current & Direct court cost & 0.1360 & [0.1359, 0.1361] \\
Council - current & Upstream cost & 0.0019 & [0.0018, 0.0020] \\
Council - current & Capacity strain & 0.0643 & [0.0637, 0.0650] \\
Council - current & Total cost index & 0.0831 & [0.0829, 0.0832] \\
Council - current & Intake acceptance & -0.0007 & [-0.0007, -0.0007] \\
Council - current & Council screening & 1.0000 & [1.0000, 1.0000] \\
Council - current & Administrative load & 0.2481 & [0.2472, 0.2491] \\
Weak form - current & Legislative response & 0.3196 & [0.3095, 0.3291] \\
Weak form - current & Response credibility & 0.0646 & [0.0631, 0.0660] \\
Weak form - current & Timely replies & 0.0608 & [0.0510, 0.0713] \\
Weak form - current & Rights index & 0.0246 & [0.0228, 0.0265] \\
Weak form - current & Compliance events & 0.3752 & [0.3592, 0.3924] \\
Weak form - current & Responsiveness & 0.0823 & [0.0798, 0.0847] \\
Weak form - current & Veto-risk index & -0.3463 & [-0.3480, -0.3446] \\
Weak form - current & Total cost index & -0.0925 & [-0.0927, -0.0924] \\
Weak form - current & Reply delay & 0.4157 & [0.4136, 0.4177] \\
Fast - reasoned & Reasons per order & -0.6954 & [-0.7258, -0.6631] \\
Fast - reasoned & Merits follow-up & -0.9917 & [-1.0000, -0.9750] \\
Fast - reasoned & Review-time input & 0.0000 & [0.0000, 0.0000] \\
Fast - reasoned & Total cost index & -0.0044 & [-0.0046, -0.0043] \\
Fast - reasoned & Emergency irregularity & 0.3081 & [0.3007, 0.3153] \\
Fast - reasoned & Delay-cost index & -0.0066 & [-0.0068, -0.0064] \\
Full court - council & Direct court cost & -0.0973 & [-0.0974, -0.0971] \\
Full court - council & Upstream cost & 0.0000 & [0.0000, 0.0001] \\
Full court - council & Capacity strain & -0.0434 & [-0.0440, -0.0428] \\
Full court - council & Total cost index & -0.0503 & [-0.0504, -0.0502] \\
Full court - council & Intake acceptance & 0.0000 & [0.0000, 0.0000] \\
Full court - council & Council screening & -1.0000 & [-1.0000, -1.0000] \\
Full court - council & Administrative load & -0.2404 & [-0.2413, -0.2395] \\
Half - weak form & Legislative response & -0.1609 & [-0.1703, -0.1521] \\
Half - weak form & Response credibility & -0.3074 & [-0.3087, -0.3061] \\
Half - weak form & Timely replies & -0.0608 & [-0.0713, -0.0510] \\
Half - weak form & Rights index & -0.0129 & [-0.0139, -0.0119] \\
Half - weak form & Compliance events & -0.0074 & [-0.0240, 0.0096] \\
Half - weak form & Responsiveness & -0.0249 & [-0.0265, -0.0234] \\
Half - weak form & Veto-risk index & 0.0258 & [0.0244, 0.0273] \\
Half - weak form & Total cost index & -0.0015 & [-0.0018, -0.0013] \\
Half - weak form & Reply delay & 0.0821 & [0.0792, 0.0849] \\
Off - weak form & Legislative response & -0.3196 & [-0.3291, -0.3095] \\
Off - weak form & Response credibility & -0.6147 & [-0.6165, -0.6130] \\
Off - weak form & Timely replies & -0.0608 & [-0.0713, -0.0510] \\
Off - weak form & Rights index & -0.0211 & [-0.0222, -0.0200] \\
Off - weak form & Compliance events & -0.0166 & [-0.0347, 0.0007] \\
Off - weak form & Responsiveness & -0.0485 & [-0.0502, -0.0469] \\
Off - weak form & Veto-risk index & 0.0510 & [0.0495, 0.0525] \\
Off - weak form & Total cost index & -0.0030 & [-0.0032, -0.0028] \\
Off - weak form & Reply delay & -0.4157 & [-0.4177, -0.4136] \\
Half - current & Legislative response & 0.1586 & [0.1515, 0.1659] \\
Half - current & Response credibility & -0.2429 & [-0.2437, -0.2421] \\
Half - current & Timely replies & 0.0000 & [0.0000, 0.0000] \\
Half - current & Rights index & 0.0117 & [0.0098, 0.0136] \\
Half - current & Compliance events & 0.3678 & [0.3497, 0.3851] \\
Half - current & Responsiveness & 0.0574 & [0.0550, 0.0598] \\
Half - current & Veto-risk index & -0.3205 & [-0.3218, -0.3193] \\
Half - current & Total cost index & -0.0941 & [-0.0943, -0.0939] \\
Half - current & Reply delay & 0.4977 & [0.4949, 0.5005] \\
Off - current & Legislative response & 0.0000 & [0.0000, 0.0000] \\
Off - current & Response credibility & -0.5502 & [-0.5510, -0.5494] \\
Off - current & Timely replies & 0.0000 & [0.0000, 0.0000] \\
Off - current & Rights index & 0.0036 & [0.0015, 0.0057] \\
Off - current & Compliance events & 0.3586 & [0.3412, 0.3748] \\
Off - current & Responsiveness & 0.0338 & [0.0316, 0.0360] \\
Off - current & Veto-risk index & -0.2953 & [-0.2958, -0.2948] \\
Off - current & Total cost index & -0.0956 & [-0.0958, -0.0954] \\
Off - current & Reply delay & 0.0000 & [0.0000, 0.0000] \\
\bottomrule
\end{longtable}
\par\noindent\begin{minipage}{0.98\linewidth}\footnotesize All 71 component/contrast combinations in this reference cell are shown without significance filtering. The CSV retains all 3,692 estimates, intervals and Monte Carlo standard errors. Intervals use 2,000 paired whole-run bootstrap draws over 120 runs; they do not measure structural or normative uncertainty.\end{minipage}
}


The review-time endpoint is a shared case input, so its paired contrasts are zero by construction and cannot support a procedural-duration claim. Conditional emergency and response rates also require care: a run without emergency orders contributes zero to the registered run-mean rate, but has no eligible order on which reasons could be observed. No legislative response is not a late or instantaneous response. Table~\ref{tab:supp-historical-denominators} preserves the low-emergency counts. Its pooled rates are explicitly post-estimate explanations, not replacement endpoints, and the sparse orders do not establish a reliable conditional comparison.

% Auto-generated by paper/scripts/generate_historical_tables.py
{\small
\setlength{\tabcolsep}{3pt}
\begin{longtable}{@{}L{0.24\linewidth}L{0.14\linewidth}rrrr@{}}
\caption{Low-emergency reason-giving denominators.}\label{tab:supp-historical-denominators}\\
\toprule
Forecast & Design & Orders & Reasons & Empty runs & Pooled rate \\
\midrule
\endfirsthead
\toprule
Forecast & Design & Orders & Reasons & Empty runs & Pooled rate \\
\midrule
\endhead
Fixed generator & Current & 583 & 13 & 0 & 0.0223 \\
Fixed generator & Reasoned & 2 & 1 & 118 & 0.5000 \\
Recency-weighted & Current & 611 & 20 & 0 & 0.0327 \\
Recency-weighted & Reasoned & 4 & 4 & 116 & 1.0000 \\
Expanding history & Current & 579 & 15 & 1 & 0.0259 \\
Expanding history & Reasoned & 1 & 1 & 119 & 1.0000 \\
Previous term & Current & 591 & 20 & 0 & 0.0338 \\
Previous term & Reasoned & 1 & 1 & 119 & 1.0000 \\
\bottomrule
\end{longtable}
\par\noindent\begin{minipage}{0.98\linewidth}\footnotesize Each row covers 120 runs. Legacy no-order run rates equal zero. Pooled reasons per order are post-estimate diagnostics, not replacement primary estimands; sparse counts cannot establish a reliable conditional comparison. An empty pooled denominator is NA, not an observed zero.\end{minipage}
}

% Auto-generated by paper/scripts/generate_historical_tables.py
{\small
\setlength{\tabcolsep}{3pt}
\begin{longtable}{@{}L{0.47\linewidth}L{0.35\linewidth}@{}}
\caption{Council-minus-current accounting-view sensitivity.}\label{tab:supp-historical-accounting}\\
\toprule
Accounting contrast & Difference range \\
\midrule
\endfirsthead
\toprule
Accounting contrast & Difference range \\
\midrule
\endhead
Direct-only view & 0.1358 to 0.1375 \\
Full accounting & 0.0828 to 0.0860 \\
Full minus direct contrast & -0.0530 to -0.0515 \\
\bottomrule
\end{longtable}
\par\noindent\begin{minipage}{0.98\linewidth}\footnotesize Direct-only versus full accounting changes the practical sign in 0 of 52 cells. These are differently constructed normalized indices, not additive monetary costs. Full per-cell paired intervals and the paired difference-of-contrasts are retained in the cost-accounting CSV.\end{minipage}
}


\subsection*{Normative sensitivity and verification boundary}

The new preference exercise samples 1,000 Dirichlet$(1,\ldots,1)$ vectors by drawing twelve positive uniforms from Python's seeded standard-library generator, applying $-\log U$, and dividing by their sum. It uses seed 2026091202, the metric order in Table~\ref{tab:supp-historical-weights}, fixed $1-x$ inversions, and no observed-grid rescaling. Scores use equal-run metric means. The same vectors apply to all designs/cells, and near-ties within $10^{-12}$ split leadership credit. Half/off competitions replace only the weak-form entry among the four base designs. Preference percentile ranges on weighted response-score changes are not Monte Carlo confidence intervals.

% Auto-generated by paper/scripts/generate_historical_tables.py
{\small
\setlength{\tabcolsep}{3pt}
\begin{longtable}{@{}L{0.44\linewidth}L{0.44\linewidth}@{}}
\caption{Metric directions in the new normative exercise.}\label{tab:supp-historical-weights}\\
\toprule
Higher-is-better & Lower-is-better \\
\midrule
\endfirsthead
\toprule
Higher-is-better & Lower-is-better \\
\midrule
\endhead
Rights index & Total cost index \\
Responsiveness & Veto-risk index \\
Legal stability & Review-time input \\
Public-legitimacy proxy & Emergency irregularity \\
Compliance events & Conflict index \\
Case-selection access & Administrative load \\
\bottomrule
\end{longtable}
\par\noindent\begin{minipage}{0.98\linewidth}\footnotesize Metric order is the higher-is-better column followed by the lower-is-better column. Only the latter is transformed to $1-x$. No observed-grid normalization is fitted. The review-time input is shared, so its weighted term cannot distinguish designs within a cell.\end{minipage}
}

% Auto-generated by paper/scripts/generate_historical_tables.py
{\small
\setlength{\tabcolsep}{3pt}
\begin{longtable}{@{}L{0.28\linewidth}L{0.32\linewidth}L{0.25\linewidth}@{}}
\caption{Normative leadership across the registered grid.}\label{tab:supp-historical-leadership}\\
\toprule
Response setting & Design & Leadership range \\
\midrule
\endfirsthead
\toprule
Response setting & Design & Leadership range \\
\midrule
\endhead
full response & Current & 0.0\% to 0.0\% \\
full response & Reasoned & 0.0\% to 63.9\% \\
full response & Council & 0.0\% to 0.0\% \\
full response & Weak form & 36.1\% to 100.0\% \\
half response & Current & 0.0\% to 0.0\% \\
half response & Reasoned & 0.0\% to 70.1\% \\
half response & Council & 0.0\% to 0.0\% \\
half response & Weak form (half) & 29.9\% to 100.0\% \\
no response & Current & 0.0\% to 0.0\% \\
no response & Reasoned & 0.0\% to 76.8\% \\
no response & Council & 0.0\% to 0.4\% \\
no response & Weak form (off) & 23.2\% to 99.8\% \\
\bottomrule
\end{longtable}
\par\noindent\begin{minipage}{0.98\linewidth}\footnotesize Minimum/maximum shares across 52 cells using the same 1,000 Dirichlet preference vectors. Each competition has four designs; only weak form is replaced. 152 of 312 base pair/cell comparisons reverse across weights. Ties within $10^{-12}$ split credit. These are illustrative preferences, not probabilities that an institution is optimal.\end{minipage}
}


The raw audit independently reconstructs 84 run measures and all 6,240 shared docket identities from 3,993,600 cases and 5,193,352 challenged objects. The separate analysis checker verifies every paired point estimate and standard error, 93 expanded-index bootstrap intervals, all weight vectors and rankings, and conditional-denominator diagnostics. Full deterministic analysis re-execution checks all saved bytes, but independent interval verification is limited to the stated 93 intervals. Table-generation hashes bind the publication to that evidence. Passing these checks establishes computational consistency, not real-world reform effects, constitutional correctness, or a winning institutional design.

\subsection*{Reproduction and artifact contract}

The retained evidence can be checked without the large raw exports using \texttt{make historical-scores-check}, \texttt{make historical-study-check}, \texttt{make historical-analysis-check}, and \texttt{make historical-analysis-validation-check}. Full institutional regeneration uses \texttt{make historical-replication}; this sequentially reruns the grid, raw audit, paired/preference analysis, and independent analysis validation. Original public SCDB archives are reacquired and hash-checked with \texttt{make historical-sources}; \texttt{make historical-benchmark-check} then rechecks the full-field revision audit and empirical reports. Do not rerun the original acquisition or forecast-freeze commands after the test has been opened.

The simulation seed is 2026091201 and the bootstrap seed is 20260912. The recorded run compiled for Java release 21 on OpenJDK 25. Shared inputs, cases, objects, segments, compositions, response/voting streams, model settings, manifests, and analysis tables have explicit composite keys in the data contract. Large raw case/object/docket/segment files and original-source caches are omitted from the anonymous package, not from the reproduction workflow. Review-time, denominator, cost-index and source-range limitations apply equally to regenerated results. In particular, the France QPC and CJEU source-range misses remain reported in the unchanged original benchmark family below.

\section*{Core Model Weights}
This table is the reader-facing bridge between the prose methods section and the implementation. It lists the fixed weights behind the major composite measures so that the appendix is not just a dump of output matrices.
% Auto-generated by paper/scripts/generate_supplement.py
{\scriptsize
\setlength{\tabcolsep}{3pt}
\begin{longtable}{@{}L{0.20\linewidth}L{0.69\linewidth}@{}}
\caption{Core model equations and fixed weights.}\label{tab:supp-model-weights}\\
\toprule
Component & Fixed construction used in the reported campaigns \\
\midrule
\endfirsthead
\toprule
Component & Fixed construction used in the reported campaigns \\
\midrule
\endhead
Shadow-docket abuse & Emergency cases start from procedure risk: fast shadow docket 0.72, reasoned emergency panel 0.34, full-court emergency 0.20, merits-follow-up design 0.16. The score adds no-merits penalty 0.26, relief-without-merits penalty 0.18, impact from salient relief or high rights threat, executive pressure 0.12, and subtracts emergency reason-giving 0.22. \\
Legal stability & $1 - 0.32$ reversal magnitude $- 0.28$ constitutional conflict $- 0.22$ shadow abuse $- 0.08$ dissent intensity $- 0.06$ concurrence fragmentation $+ 0.12$ scenario stability preference, truncated to $[0,1]$. \\
Legitimacy & Weighted average of case public trust 0.24, dynamic public trust 0.12, transparency 0.18, majority share 0.12, rights protection 0.18, inverse partisan alignment 0.14, inverse shadow abuse 0.14, and constitutional conflict penalty 0.18. \\
Rights protection & Merits invalidation begins at 0.44 plus rights-threat weight 0.54 and legal-ambiguity penalty 0.08. Emergency relief without merits begins at 0.38 plus rights-threat weight 0.40 and time-to-review penalty 0.08. No-relief outcomes fall with rights threat and delay. Overrides reduce protection when rights threat is high. \\
Constitutional conflict & Adds countermajoritarian clash from merits invalidation or emergency relief, override penalty 0.36, shadow-abuse weight 0.30, cross-check conflict 0.24, executive pressure 0.10, lower-court panel skew 0.05, state-federal tension 0.12, and intercourt conflict 0.10. \\
Compliance and enforcement & Compliance starts at 0.46 and moves with legitimacy 0.24, dynamic compliance norm 0.22, court independence 0.06, conflict penalties, emergency-without-merits penalty 0.08, override penalty 0.10, and executive pressure 0.08. Executive implementation, agency nonacquiescence, reenactment, and local compliance are separate downstream Bernoulli draws. \\
Institutional cost & Total cost is $0.25$ direct court cost $+ 0.20$ upstream screening $+ 0.20$ capacity strain $+ 0.18$ delay $+ 0.17$ implementation complexity. The subcomponents depend on court size, review structure, docket control, emergency procedure, council screening, cross-checking, and panel/en banc design. \\
Directional score & Unweighted average of 29 display metrics after reversing lower-is-better measures. It is a diagnostic display score and not a welfare function or fitted objective. \\
Legal-transplant feasibility & Scenario transplant score 0.52, political-culture fit 0.22, case public trust 0.10, inverse institutional demand 0.10, and inverse political mismatch 0.06. Institutional demand combines total cost, upstream screening, complexity, intake screening, and a synthetic-mechanism penalty. \\
Political-culture sensitivity & Mechanism fragility baseline plus low culture fit, legislative conflict, partisan pressure, constitutional conflict, implementation complexity, and state-federal tension. Lower is better. \\
Veto-relocation risk & Low-rights interventions, council or cross-check veto points, intake-screening intensity, upstream cost, low acceptance rate, high independence, and low transparency increase risk. Legislative response and rights-threat justification reduce it. \\
Democratic constitutionalism & Legal stability 0.13, rights protection 0.20, legitimacy 0.14, democratic responsiveness 0.18, compliance 0.11, inverse conflict 0.09, inverse veto relocation 0.07, transplant feasibility 0.05, and inverse culture sensitivity 0.03. \\
\end{longtable}
}


\section*{Model-to-Code Crosswalk}
This table maps prose components to implementation modules. It is included to make the model auditable at review stage without forcing readers to infer which code path generated each output.
% Auto-generated by paper/scripts/generate_supplement.py
{\scriptsize
\setlength{\tabcolsep}{2pt}
\begin{longtable}{@{}L{0.15\linewidth}L{0.19\linewidth}L{0.23\linewidth}L{0.23\linewidth}L{0.10\linewidth}@{}}
\caption{Model-to-code crosswalk.}\label{tab:supp-model-crosswalk}\\
\toprule
Stage & Main module & Inputs & Outputs & Evidence status \\
\midrule
\endfirsthead
\toprule
Stage & Main module & Inputs & Outputs & Evidence status \\
\midrule
\endhead
Scenario catalog & ScenarioCatalog;\newline Scenario & scenario key, design configuration, world spec & review process factory and scenario metadata & institutional mapping \\
Docket generation & WorldGenerator;\newline CaseFile & world variables, imported legislative signals, random seed & case-level facts, doctrine, policy domain, lower-court path & synthetic input generator \\
Admission and access & IntakeModel;\newline IntakeEstimate & case salience, rights threat, litigant capacity, public-interest support, repeat-player advantage, docket control & reviewed flag, intake filings, screened filings, intake acceptance, case-selection access & diagnostic plus source-range proxy rows \\
Emergency docket & Emergency\hspace{0pt}Docket\hspace{0pt}Model & urgency, executive pressure, vote share, emergency procedure & relief, reasons, vote disclosure, public disagreement, merits follow-up, emergency-process irregularity & synthetic mechanism with source diagnostics \\
Merits voting & Constitutional\hspace{0pt}Review\hspace{0pt}Process;\newline VoteResult & justice traits, policy position, rights threat, mandate, lower-court signals, thresholds & strike vote share, invalidation, concurrence fragmentation, dissent intensity & stylized decision model \\
Noncourt mechanisms & Decision\hspace{0pt}Mechanism\hspace{0pt}Model;\newline Post\hspace{0pt}Decision\hspace{0pt}Response\hspace{0pt}Model & review timing, ombudsman, public defender, rights statement, declaration, override, party context & legislative response, response delay, response credibility, supranational route flags & synthetic mechanism with source diagnostics \\
Compliance and enforcement & ComplianceModel;\newline ReactionState & legitimacy proxy, constitutional-conflict index, implementation capacity, dynamic reaction state & compliance, defiance, workaround, repeated litigation, executive implementation, nonacquiescence, reenactment, local compliance & synthetic mechanism \\
Institutional cost & Institutional\hspace{0pt}Cost\hspace{0pt}Model;\newline CostEstimate & court size, structure, screening, emergency procedure, recusal, intake pressure & direct court cost, upstream screening, capacity strain, delay, complexity, total cost & normalized diagnostic \\
Outcome scoring & Outcome\hspace{0pt}Score\hspace{0pt}Model;\newline Directional\hspace{0pt}Score\hspace{0pt}Model & case outcomes, costs, access, response credibility, transplant and culture modifiers & rights protection, conflict, democratic constitutionalism, veto relocation, directional display score & synthetic diagnostic \\
Campaign and calibration & CampaignRunner;\newline Campaign\hspace{0pt}Calibration\hspace{0pt}Evaluator & scenario rows, case outcomes, calibration source-observation matrix & CSV reports, uncertainty bands, external diagnostic comparisons, provenance manifests & reporting and source-range diagnostics \\
\bottomrule
\end{longtable}
\noindent\footnotesize The crosswalk identifies where each prose model component is implemented. Evidence status separates institutional mappings and source-range diagnostics from synthetic mechanisms and normalized display scores.
}


\section*{Formal Model Appendix}

This appendix states the model in compact notation so the simulation can be evaluated without first reading the implementation source. The implementation remains deterministic conditional on the random seed, scenario configuration, world specification, and imported legislative signals.

\subsection*{Units and Inputs}

Each simulated observation is a case $i$ in review period $t$ under institutional design $d$. The generated case vector is
\[
x_i = (r_i, u_i, a_i, s_i, p_i, m_i, e_i, \ell_i, q_i, z_i),
\]
where $r_i$ is rights threat, $u_i$ urgency, $a_i$ legal ambiguity, $s_i$ constitutional salience, $p_i$ public support, $m_i$ legislative mandate, $e_i$ executive pressure, $\ell_i$ lower-court conflict, $q_i$ state-federal tension, and $z_i$ policy-domain position. A world context $w$ supplies public trust, legislative conflict, partisan pressure, party fragmentation, government control, electoral time pressure, civil-society capacity, implementation capacity, and legal-tradition compatibility. For validation presets, those context inputs are read from \texttt{config/context/country-year-context.csv}; for synthetic campaigns, they are stress-case parameters.

\subsection*{Intake and Case Selection}

The intake universe expands each generated case into expected filings $F_i$ and screened filings $S_i$:
\[
F_i = 1 + \lfloor 7A_d + 3r_i + 2s_i + 2\ell_i + H_i \rfloor,
\quad
S_i = \max(1, \lfloor F_i \alpha_i \rfloor),
\]
where $A_d$ is the scenario's structural access level, $H_i$ is an emergency or high-salience increment, and $\alpha_i$ is the intake acceptance probability after docket-control screening. Review occurs when a weighted access score exceeds the design-specific threshold:
\[
\Pr(R_i = 1) = g(A_d, r_i, s_i, a_i, \ell_i, u_i, \text{ombudsman}_i, \text{public-defender}_i, \text{rights-statement}_i).
\]
The reported \texttt{reviewRate} is $\sum R_i/N$ and \texttt{intakeAcceptanceRate} is $\sum R_i/\sum F_i$.

\subsection*{Merits and Emergency Decisions}

The primary vote share for invalidation is generated from justice ideology, rights sensitivity, institutionalism, case facts, and design thresholds. In compact form:
\[
V_i = h(J_{dt}, r_i, a_i, s_i, e_i, z_i, \tau_d, \rho_d),
\]
where $J_{dt}$ is the period-specific court composition, $\tau_d$ is the invalidation threshold, and $\rho_d$ captures recusal, panel, en banc, cross-checking, and council-screening rules. Merits invalidation occurs when $V_i \geq \tau_d$ and the case has reached merits review. Emergency relief is modeled separately from merits invalidation:
\[
E_i = 1[u_i + \eta_w > \theta_d], \quad
\text{Relief}_i = k(E_i, r_i, u_i, e_i, \text{procedure}_d, \text{reasons}_i, \text{vote-disclosure}_i).
\]
This separation is why emergency relief, reason-giving, vote disclosure, public disagreement, merits review, and merits invalidation are distinct report columns.

\subsection*{Mechanism Flags}

Design-specific mechanisms are observed rather than folded into a single court-variant label. The model records weak-form declarations, suspended declarations, formal overrides, rights-impact statements, ombudsman-triggered access, constitutional public-defender participation, pre-enactment review, abstract review, panel/en banc review, cross-checking, and constitutional-council screening. For supranational review, the CJEU-style route vector is
\[
(P_i,A_i,D_i,O_i) \sim \text{Categorical}(\pi_P(x_i), \pi_A(x_i), \pi_D(x_i), 1-\pi_P-\pi_A-\pi_D),
\]
with preliminary-reference, appeal, and direct-action probabilities bounded around the official 2024 route mix used in calibration diagnostics.

\subsection*{Legislative Response}

A legislative-response opportunity is triggered by weak-form declarations, suspended declarations, mandatory response cycles, formal override structures, or high-salience invalidations. Response credibility is
\[
\begin{aligned}
C_i={}&c_0+c_1r_i+c_2p_i+c_3m_i
+c_4\text{civil-society}_w+c_5\text{implementation}_w\\
&-c_6\text{fragmentation}_w-c_7\text{electoral-pressure}_w.
\end{aligned}
\]
The reported response rate is the share of cases with a response; delay is normalized to the $[0,1]$ review-period scale; timely response is the share of responses completed before the scenario deadline.

\subsection*{Compliance and Enforcement}

Compliance is drawn after the decision and political reaction update:
\[
\begin{aligned}
\Pr(\text{comply}_i)={}&b_0+b_1\text{legitimacy}_i
+b_2\text{implementation}_w+b_3\text{independence}_d\\
&-b_4\text{conflict}_t-b_5e_i-b_6\text{shadow-abuse}_i.
\end{aligned}
\]
The simulator separately records executive implementation, agency nonacquiescence, legislative reenactment, local-government compliance, defiance, workaround behavior, and repeated litigation. These channels prevent a high rights-protection score from being treated as automatically implemented law.

\subsection*{Composite Diagnostics}

The principal composite diagnostic is democratic constitutionalism:
\[
\begin{aligned}
DC_i ={}& .11\,LS_i + .18\,RP_i + .12\,LG_i + .15\,DR_i + .09\,CP_i
+ .08(1-CC_i) + .07(1-VR_i) \\
&+ .05\,TF_i + .03(1-PC_i) + .03\,LC_i + .03\,CA_i
+ .03(1-GR_i) + .03\,IC_i.
\end{aligned}
\]
Here $LS$ is legal stability, $RP$ rights protection, $LG$ legitimacy, $DR$ democratic responsiveness, $CP$ compliance, $CC$ constitutional conflict, $VR$ veto-relocation risk, $TF$ transplant feasibility, $PC$ political-culture sensitivity, $LC$ legislative-response credibility, $CA$ case-selection access, $GR$ government repeat-player advantage, and $IC$ implementation capacity. The directional score shown in campaign tables is an unweighted display average after reversing lower-is-better metrics. It is not a fitted welfare function.

\subsection*{Uncertainty and Calibration}

Campaign uncertainty bands use a cluster bootstrap over run-level outputs. Validation-style diagnostics compare simulated columns to source-backed calibration rows only when a row has a public URL, nonzero denominator, construction note, and direct simulator analogue. Context rows, public-trust proxies, normalized cost estimates, and unverified source candidates remain visible in the supplement but are not counted as empirical validation evidence.


\section*{Starter Comparison Matrix}
The main comparison matrix is a complete audit trail for the baseline and stress-case rows summarized in the main manuscript. It shows how the same institutional designs move across shared docket stresses. Mechanism labels make clear which rows are court-architecture comparisons and which rows are devices such as weak-form review, rights-impact statements, or mandatory legislative response cycles.
% Auto-generated by paper/scripts/generate_supplement.py
{\scriptsize
\setlength{\tabcolsep}{2pt}
\begin{longtable}{@{}L{0.12\linewidth}L{0.23\linewidth}L{0.11\linewidth}rrrrrr@{}}
\caption{Full main-campaign scenario matrix.}\label{tab:supp-scenario-matrix}\\
\toprule
Case & Scenario & Mechanism & Score & Dem. const. & Rights & Veto $\downarrow$ & Culture $\downarrow$ & Cost $\downarrow$ \\
\midrule
\endfirsthead
\toprule
Case & Scenario & Mechanism & Score & Dem. const. & Rights & Veto $\downarrow$ & Culture $\downarrow$ & Cost $\downarrow$ \\
\midrule
\endhead
Baseline & Current-style federal supreme court & Strong form & 0.525 & 0.571 & 0.687 & 0.529 & 0.728 & 0.450 \\
Baseline & Staggered 18-year nonrenewable terms & Strong form & 0.535 & 0.578 & 0.695 & 0.520 & 0.727 & 0.455 \\
Baseline & Fifteen-justice supermajority commission court & Strong form & 0.615 & 0.643 & 0.739 & 0.436 & 0.717 & 0.470 \\
Baseline & Supermajority required to invalidate laws & Strong form & 0.604 & 0.638 & 0.736 & 0.446 & 0.720 & 0.457 \\
Baseline & Strict recusal with substitute justices & Strong form & 0.622 & 0.648 & 0.739 & 0.434 & 0.720 & 0.476 \\
Baseline & Reasoned emergency orders with merits follow-up & Strong form & 0.625 & 0.649 & 0.740 & 0.446 & 0.719 & 0.460 \\
Baseline & Panel review with en banc safety valve & Strong form & 0.609 & 0.643 & 0.739 & 0.433 & 0.720 & 0.511 \\
Baseline & Dual supreme courts with cross-checking invalidations & Strong form & 0.628 & 0.650 & 0.747 & 0.596 & 0.733 & 0.589 \\
Baseline & Pre-enactment constitutional council & Strong form & 0.632 & 0.652 & 0.741 & 0.662 & 0.726 & 0.533 \\
Baseline & Judicial review with delayed legislative override & Strong form & 0.592 & 0.637 & 0.733 & 0.435 & 0.724 & 0.466 \\
Baseline & Retention-election accountability court & Strong form & 0.622 & 0.656 & 0.729 & 0.407 & 0.725 & 0.482 \\
Baseline & Hybrid court balancing independence and accountability & Strong form & 0.640 & 0.663 & 0.746 & 0.433 & 0.732 & 0.542 \\
Baseline & Weak-form review with legislative reply & Weak form & 0.699 & 0.709 & 0.718 & 0.177 & 0.759 & 0.359 \\
Baseline & Suspended declarations of invalidity & Suspended & 0.675 & 0.680 & 0.724 & 0.207 & 0.836 & 0.383 \\
Baseline & Strong-form review with explicit override clause & Override & 0.614 & 0.655 & 0.738 & 0.355 & 0.862 & 0.425 \\
Baseline & Pre-enactment review before laws take effect & Pre-enact. & 0.652 & 0.667 & 0.766 & 0.511 & 0.833 & 0.437 \\
Baseline & Abstract review tribunal & Abstract & 0.643 & 0.665 & 0.770 & 0.312 & 0.905 & 0.436 \\
Baseline & Ombudsman-triggered constitutional review & Ombudsman & 0.607 & 0.632 & 0.736 & 0.525 & 0.775 & 0.516 \\
Baseline & Constitutional public defender access model & Defender & 0.618 & 0.646 & 0.773 & 0.514 & 0.746 & 0.518 \\
Baseline & Rights-impact statements before review & Rights stmt. & 0.688 & 0.691 & 0.805 & 0.502 & 0.709 & 0.437 \\
Baseline & Mandatory legislative response cycles & Response & 0.659 & 0.680 & 0.717 & 0.224 & 0.916 & 0.394 \\
Polarized appointments & Current-style federal supreme court & Strong form & 0.506 & 0.550 & 0.671 & 0.532 & 0.767 & 0.450 \\
Polarized appointments & Staggered 18-year nonrenewable terms & Strong form & 0.515 & 0.556 & 0.678 & 0.523 & 0.767 & 0.455 \\
Polarized appointments & Fifteen-justice supermajority commission court & Strong form & 0.594 & 0.624 & 0.735 & 0.439 & 0.755 & 0.470 \\
Polarized appointments & Supermajority required to invalidate laws & Strong form & 0.582 & 0.618 & 0.729 & 0.449 & 0.759 & 0.457 \\
Polarized appointments & Strict recusal with substitute justices & Strong form & 0.601 & 0.629 & 0.735 & 0.437 & 0.759 & 0.476 \\
Polarized appointments & Reasoned emergency orders with merits follow-up & Strong form & 0.602 & 0.628 & 0.733 & 0.450 & 0.758 & 0.460 \\
Polarized appointments & Panel review with en banc safety valve & Strong form & 0.592 & 0.626 & 0.735 & 0.436 & 0.758 & 0.512 \\
Polarized appointments & Dual supreme courts with cross-checking invalidations & Strong form & 0.604 & 0.630 & 0.744 & 0.598 & 0.772 & 0.589 \\
Polarized appointments & Pre-enactment constitutional council & Strong form & 0.609 & 0.631 & 0.739 & 0.664 & 0.765 & 0.533 \\
Polarized appointments & Judicial review with delayed legislative override & Strong form & 0.570 & 0.616 & 0.726 & 0.438 & 0.763 & 0.466 \\
Polarized appointments & Retention-election accountability court & Strong form & 0.593 & 0.631 & 0.725 & 0.408 & 0.764 & 0.482 \\
Polarized appointments & Hybrid court balancing independence and accountability & Strong form & 0.617 & 0.642 & 0.743 & 0.435 & 0.771 & 0.543 \\
Polarized appointments & Weak-form review with legislative reply & Weak form & 0.686 & 0.693 & 0.712 & 0.173 & 0.797 & 0.359 \\
Polarized appointments & Suspended declarations of invalidity & Suspended & 0.655 & 0.663 & 0.716 & 0.205 & 0.875 & 0.383 \\
Polarized appointments & Strong-form review with explicit override clause & Override & 0.587 & 0.631 & 0.731 & 0.351 & 0.900 & 0.426 \\
Polarized appointments & Pre-enactment review before laws take effect & Pre-enact. & 0.631 & 0.647 & 0.766 & 0.513 & 0.871 & 0.438 \\
Polarized appointments & Abstract review tribunal & Abstract & 0.622 & 0.647 & 0.769 & 0.314 & 0.943 & 0.436 \\
Polarized appointments & Ombudsman-triggered constitutional review & Ombudsman & 0.592 & 0.616 & 0.733 & 0.527 & 0.814 & 0.517 \\
Polarized appointments & Constitutional public defender access model & Defender & 0.604 & 0.631 & 0.771 & 0.516 & 0.784 & 0.518 \\
Polarized appointments & Rights-impact statements before review & Rights stmt. & 0.662 & 0.668 & 0.802 & 0.504 & 0.747 & 0.437 \\
Polarized appointments & Mandatory legislative response cycles & Response & 0.636 & 0.658 & 0.712 & 0.221 & 0.954 & 0.395 \\
Rights-threat surge & Current-style federal supreme court & Strong form & 0.497 & 0.555 & 0.679 & 0.504 & 0.763 & 0.452 \\
Rights-threat surge & Staggered 18-year nonrenewable terms & Strong form & 0.504 & 0.558 & 0.680 & 0.495 & 0.763 & 0.457 \\
Rights-threat surge & Fifteen-justice supermajority commission court & Strong form & 0.569 & 0.620 & 0.731 & 0.411 & 0.753 & 0.472 \\
Rights-threat surge & Supermajority required to invalidate laws & Strong form & 0.565 & 0.617 & 0.728 & 0.421 & 0.756 & 0.459 \\
Rights-threat surge & Strict recusal with substitute justices & Strong form & 0.575 & 0.622 & 0.730 & 0.409 & 0.756 & 0.478 \\
Rights-threat surge & Reasoned emergency orders with merits follow-up & Strong form & 0.584 & 0.631 & 0.742 & 0.419 & 0.754 & 0.462 \\
Rights-threat surge & Panel review with en banc safety valve & Strong form & 0.566 & 0.621 & 0.732 & 0.408 & 0.755 & 0.516 \\
Rights-threat surge & Dual supreme courts with cross-checking invalidations & Strong form & 0.572 & 0.620 & 0.737 & 0.572 & 0.769 & 0.591 \\
Rights-threat surge & Pre-enactment constitutional council & Strong form & 0.573 & 0.620 & 0.725 & 0.639 & 0.762 & 0.535 \\
Rights-threat surge & Judicial review with delayed legislative override & Strong form & 0.550 & 0.616 & 0.727 & 0.409 & 0.759 & 0.468 \\
Rights-threat surge & Retention-election accountability court & Strong form & 0.558 & 0.617 & 0.702 & 0.380 & 0.761 & 0.484 \\
Rights-threat surge & Hybrid court balancing independence and accountability & Strong form & 0.584 & 0.635 & 0.738 & 0.408 & 0.768 & 0.548 \\
Rights-threat surge & Weak-form review with legislative reply & Weak form & 0.672 & 0.684 & 0.677 & 0.106 & 0.792 & 0.365 \\
Rights-threat surge & Suspended declarations of invalidity & Suspended & 0.631 & 0.660 & 0.686 & 0.136 & 0.874 & 0.389 \\
Rights-threat surge & Strong-form review with explicit override clause & Override & 0.566 & 0.629 & 0.730 & 0.320 & 0.897 & 0.428 \\
Rights-threat surge & Pre-enactment review before laws take effect & Pre-enact. & 0.592 & 0.634 & 0.754 & 0.490 & 0.869 & 0.441 \\
Rights-threat surge & Abstract review tribunal & Abstract & 0.594 & 0.642 & 0.772 & 0.287 & 0.941 & 0.443 \\
Rights-threat surge & Ombudsman-triggered constitutional review & Ombudsman & 0.563 & 0.612 & 0.734 & 0.503 & 0.812 & 0.524 \\
Rights-threat surge & Constitutional public defender access model & Defender & 0.578 & 0.638 & 0.823 & 0.489 & 0.783 & 0.526 \\
Rights-threat surge & Rights-impact statements before review & Rights stmt. & 0.621 & 0.655 & 0.792 & 0.480 & 0.745 & 0.440 \\
Rights-threat surge & Mandatory legislative response cycles & Response & 0.610 & 0.653 & 0.689 & 0.153 & 0.948 & 0.400 \\
Emergency docket stress & Current-style federal supreme court & Strong form & 0.485 & 0.484 & 0.607 & 0.524 & 0.774 & 0.455 \\
Emergency docket stress & Staggered 18-year nonrenewable terms & Strong form & 0.494 & 0.493 & 0.622 & 0.516 & 0.773 & 0.460 \\
Emergency docket stress & Fifteen-justice supermajority commission court & Strong form & 0.557 & 0.583 & 0.700 & 0.434 & 0.755 & 0.479 \\
Emergency docket stress & Supermajority required to invalidate laws & Strong form & 0.545 & 0.573 & 0.689 & 0.443 & 0.760 & 0.466 \\
Emergency docket stress & Strict recusal with substitute justices & Strong form & 0.565 & 0.587 & 0.701 & 0.432 & 0.758 & 0.485 \\
Emergency docket stress & Reasoned emergency orders with merits follow-up & Strong form & 0.598 & 0.627 & 0.729 & 0.447 & 0.745 & 0.470 \\
Emergency docket stress & Panel review with en banc safety valve & Strong form & 0.555 & 0.585 & 0.701 & 0.430 & 0.757 & 0.521 \\
Emergency docket stress & Dual supreme courts with cross-checking invalidations & Strong form & 0.593 & 0.620 & 0.734 & 0.596 & 0.760 & 0.599 \\
Emergency docket stress & Pre-enactment constitutional council & Strong form & 0.563 & 0.582 & 0.704 & 0.661 & 0.764 & 0.542 \\
Emergency docket stress & Judicial review with delayed legislative override & Strong form & 0.537 & 0.576 & 0.692 & 0.433 & 0.762 & 0.475 \\
Emergency docket stress & Retention-election accountability court & Strong form & 0.556 & 0.595 & 0.710 & 0.408 & 0.762 & 0.491 \\
Emergency docket stress & Hybrid court balancing independence and accountability & Strong form & 0.589 & 0.636 & 0.734 & 0.433 & 0.758 & 0.553 \\
Emergency docket stress & Weak-form review with legislative reply & Weak form & 0.608 & 0.616 & 0.671 & 0.227 & 0.798 & 0.376 \\
Emergency docket stress & Suspended declarations of invalidity & Suspended & 0.643 & 0.657 & 0.711 & 0.245 & 0.863 & 0.403 \\
Emergency docket stress & Strong-form review with explicit override clause & Override & 0.546 & 0.586 & 0.699 & 0.373 & 0.902 & 0.434 \\
Emergency docket stress & Pre-enactment review before laws take effect & Pre-enact. & 0.581 & 0.594 & 0.735 & 0.550 & 0.872 & 0.457 \\
Emergency docket stress & Abstract review tribunal & Abstract & 0.609 & 0.633 & 0.759 & 0.356 & 0.932 & 0.458 \\
Emergency docket stress & Ombudsman-triggered constitutional review & Ombudsman & 0.554 & 0.573 & 0.695 & 0.522 & 0.813 & 0.526 \\
Emergency docket stress & Constitutional public defender access model & Defender & 0.596 & 0.628 & 0.767 & 0.514 & 0.772 & 0.528 \\
Emergency docket stress & Rights-impact statements before review & Rights stmt. & 0.611 & 0.615 & 0.776 & 0.540 & 0.748 & 0.456 \\
Emergency docket stress & Mandatory legislative response cycles & Response & 0.624 & 0.652 & 0.715 & 0.263 & 0.941 & 0.415 \\
Low-trust conflict & Current-style federal supreme court & Strong form & 0.504 & 0.545 & 0.676 & 0.534 & 0.800 & 0.456 \\
Low-trust conflict & Staggered 18-year nonrenewable terms & Strong form & 0.512 & 0.552 & 0.685 & 0.525 & 0.799 & 0.461 \\
Low-trust conflict & Fifteen-justice supermajority commission court & Strong form & 0.581 & 0.609 & 0.732 & 0.442 & 0.789 & 0.476 \\
Low-trust conflict & Supermajority required to invalidate laws & Strong form & 0.576 & 0.607 & 0.730 & 0.451 & 0.793 & 0.463 \\
Low-trust conflict & Strict recusal with substitute justices & Strong form & 0.586 & 0.612 & 0.732 & 0.440 & 0.793 & 0.482 \\
Low-trust conflict & Reasoned emergency orders with merits follow-up & Strong form & 0.594 & 0.616 & 0.733 & 0.452 & 0.791 & 0.466 \\
Low-trust conflict & Panel review with en banc safety valve & Strong form & 0.576 & 0.609 & 0.731 & 0.439 & 0.792 & 0.521 \\
Low-trust conflict & Dual supreme courts with cross-checking invalidations & Strong form & 0.586 & 0.610 & 0.742 & 0.601 & 0.806 & 0.595 \\
Low-trust conflict & Pre-enactment constitutional council & Strong form & 0.585 & 0.608 & 0.736 & 0.667 & 0.799 & 0.539 \\
Low-trust conflict & Judicial review with delayed legislative override & Strong form & 0.561 & 0.604 & 0.725 & 0.438 & 0.797 & 0.472 \\
Low-trust conflict & Retention-election accountability court & Strong form & 0.573 & 0.612 & 0.721 & 0.409 & 0.798 & 0.487 \\
Low-trust conflict & Hybrid court balancing independence and accountability & Strong form & 0.599 & 0.624 & 0.740 & 0.437 & 0.805 & 0.552 \\
Low-trust conflict & Weak-form review with legislative reply & Weak form & 0.645 & 0.657 & 0.709 & 0.177 & 0.830 & 0.366 \\
Low-trust conflict & Suspended declarations of invalidity & Suspended & 0.628 & 0.637 & 0.687 & 0.208 & 0.909 & 0.389 \\
Low-trust conflict & Strong-form review with explicit override clause & Override & 0.569 & 0.614 & 0.731 & 0.375 & 0.935 & 0.430 \\
Low-trust conflict & Pre-enactment review before laws take effect & Pre-enact. & 0.602 & 0.621 & 0.754 & 0.516 & 0.905 & 0.443 \\
Low-trust conflict & Abstract review tribunal & Abstract & 0.603 & 0.628 & 0.757 & 0.316 & 0.976 & 0.445 \\
Low-trust conflict & Ombudsman-triggered constitutional review & Ombudsman & 0.578 & 0.600 & 0.728 & 0.531 & 0.848 & 0.522 \\
Low-trust conflict & Constitutional public defender access model & Defender & 0.588 & 0.614 & 0.768 & 0.520 & 0.819 & 0.524 \\
Low-trust conflict & Rights-impact statements before review & Rights stmt. & 0.635 & 0.643 & 0.801 & 0.507 & 0.781 & 0.444 \\
Low-trust conflict & Mandatory legislative response cycles & Response & 0.612 & 0.635 & 0.705 & 0.224 & 0.985 & 0.401 \\
\bottomrule
\end{longtable}
}


\section*{Sensitivity Matrix}
The sensitivity matrix reports the same institutional designs under high and low assumptions for emergency pressure, appointment polarization, rights-threat rate, public trust, and legislature-court conflict. The purpose is to expose parameter dependence before treating any design comparison as substantively stable. The columns emphasize democratic constitutionalism, veto-relocation risk, transplant feasibility, and political-culture sensitivity because those are the key additions for evaluating review as democratic infrastructure rather than as judicial veto power.
% Auto-generated by paper/scripts/generate_supplement.py
{\scriptsize
\setlength{\tabcolsep}{2pt}
\begin{longtable}{@{}L{0.19\linewidth}L{0.22\linewidth}rrrrrr@{}}
\caption{Sensitivity-campaign results matrix.}\label{tab:supp-sensitivity}\\
\toprule
Stress case & Scenario & Score & Dem. const. & Veto $\downarrow$ & Transplant & Culture $\downarrow$ & Trust \\
\midrule
\endfirsthead
\toprule
Stress case & Scenario & Score & Dem. const. & Veto $\downarrow$ & Transplant & Culture $\downarrow$ & Trust \\
\midrule
\endhead
Baseline & Current-style federal supreme court & 0.526 & 0.572 & 0.529 & 0.644 & 0.728 & 0.028 \\
Baseline & Staggered 18-year nonrenewable terms & 0.535 & 0.578 & 0.520 & 0.645 & 0.727 & 0.029 \\
Baseline & Fifteen-justice supermajority commission court & 0.615 & 0.643 & 0.436 & 0.645 & 0.717 & 0.056 \\
Baseline & Supermajority required to invalidate laws & 0.604 & 0.638 & 0.446 & 0.644 & 0.720 & 0.050 \\
Baseline & Strict recusal with substitute justices & 0.623 & 0.647 & 0.434 & 0.644 & 0.720 & 0.060 \\
Baseline & Reasoned emergency orders with merits follow-up & 0.624 & 0.649 & 0.446 & 0.644 & 0.719 & 0.060 \\
Baseline & Panel review with en banc safety valve & 0.608 & 0.643 & 0.433 & 0.644 & 0.719 & 0.052 \\
Baseline & Dual supreme courts with cross-checking invalidations & 0.629 & 0.650 & 0.596 & 0.641 & 0.733 & 0.071 \\
Baseline & Pre-enactment constitutional council & 0.631 & 0.652 & 0.662 & 0.643 & 0.726 & 0.079 \\
Baseline & Judicial review with delayed legislative override & 0.593 & 0.637 & 0.434 & 0.644 & 0.724 & 0.048 \\
Baseline & Retention-election accountability court & 0.622 & 0.656 & 0.407 & 0.644 & 0.725 & 0.068 \\
Baseline & Hybrid court balancing independence and accountability & 0.640 & 0.662 & 0.432 & 0.642 & 0.732 & 0.073 \\
Baseline & Weak-form review with legislative reply & 0.699 & 0.709 & 0.177 & 0.678 & 0.759 & 0.196 \\
Baseline & Suspended declarations of invalidity & 0.673 & 0.680 & 0.207 & 0.624 & 0.836 & 0.062 \\
Baseline & Strong-form review with explicit override clause & 0.616 & 0.655 & 0.355 & 0.632 & 0.862 & 0.070 \\
Baseline & Pre-enactment review before laws take effect & 0.653 & 0.667 & 0.511 & 0.581 & 0.833 & 0.084 \\
Baseline & Abstract review tribunal & 0.645 & 0.665 & 0.311 & 0.565 & 0.905 & 0.068 \\
Baseline & Ombudsman-triggered constitutional review & 0.605 & 0.631 & 0.525 & 0.604 & 0.775 & 0.051 \\
Baseline & Constitutional public defender access model & 0.618 & 0.645 & 0.513 & 0.632 & 0.746 & 0.056 \\
Baseline & Rights-impact statements before review & 0.688 & 0.691 & 0.502 & 0.699 & 0.709 & 0.098 \\
Baseline & Mandatory legislative response cycles & 0.659 & 0.680 & 0.224 & 0.616 & 0.916 & 0.089 \\
Low emergency pressure & Current-style federal supreme court & 0.549 & 0.612 & 0.532 & 0.644 & 0.717 & 0.041 \\
Low emergency pressure & Staggered 18-year nonrenewable terms & 0.556 & 0.617 & 0.522 & 0.645 & 0.716 & 0.042 \\
Low emergency pressure & Fifteen-justice supermajority commission court & 0.634 & 0.655 & 0.436 & 0.645 & 0.714 & 0.067 \\
Low emergency pressure & Supermajority required to invalidate laws & 0.619 & 0.651 & 0.446 & 0.644 & 0.718 & 0.063 \\
Low emergency pressure & Strict recusal with substitute justices & 0.636 & 0.660 & 0.434 & 0.644 & 0.717 & 0.074 \\
Low emergency pressure & Reasoned emergency orders with merits follow-up & 0.636 & 0.651 & 0.446 & 0.644 & 0.718 & 0.062 \\
Low emergency pressure & Panel review with en banc safety valve & 0.626 & 0.655 & 0.434 & 0.644 & 0.717 & 0.065 \\
Low emergency pressure & Dual supreme courts with cross-checking invalidations & 0.633 & 0.653 & 0.596 & 0.641 & 0.733 & 0.073 \\
Low emergency pressure & Pre-enactment constitutional council & 0.648 & 0.666 & 0.662 & 0.643 & 0.724 & 0.105 \\
Low emergency pressure & Judicial review with delayed legislative override & 0.607 & 0.648 & 0.435 & 0.644 & 0.721 & 0.057 \\
Low emergency pressure & Retention-election accountability court & 0.635 & 0.667 & 0.406 & 0.644 & 0.722 & 0.085 \\
Low emergency pressure & Hybrid court balancing independence and accountability & 0.643 & 0.666 & 0.432 & 0.642 & 0.732 & 0.080 \\
Low emergency pressure & Weak-form review with legislative reply & 0.711 & 0.724 & 0.166 & 0.678 & 0.756 & 0.447 \\
Low emergency pressure & Suspended declarations of invalidity & 0.687 & 0.685 & 0.202 & 0.624 & 0.836 & 0.069 \\
Low emergency pressure & Strong-form review with explicit override clause & 0.634 & 0.671 & 0.352 & 0.632 & 0.859 & 0.091 \\
Low emergency pressure & Pre-enactment review before laws take effect & 0.668 & 0.682 & 0.504 & 0.581 & 0.830 & 0.112 \\
Low emergency pressure & Abstract review tribunal & 0.650 & 0.669 & 0.307 & 0.565 & 0.904 & 0.075 \\
Low emergency pressure & Ombudsman-triggered constitutional review & 0.621 & 0.645 & 0.526 & 0.604 & 0.773 & 0.068 \\
Low emergency pressure & Constitutional public defender access model & 0.621 & 0.649 & 0.513 & 0.632 & 0.745 & 0.061 \\
Low emergency pressure & Rights-impact statements before review & 0.707 & 0.704 & 0.494 & 0.699 & 0.706 & 0.147 \\
Low emergency pressure & Mandatory legislative response cycles & 0.671 & 0.686 & 0.219 & 0.616 & 0.915 & 0.107 \\
High emergency pressure & Current-style federal supreme court & 0.490 & 0.488 & 0.524 & 0.644 & 0.753 & 0.014 \\
High emergency pressure & Staggered 18-year nonrenewable terms & 0.498 & 0.497 & 0.516 & 0.645 & 0.752 & 0.015 \\
High emergency pressure & Fifteen-justice supermajority commission court & 0.561 & 0.586 & 0.433 & 0.644 & 0.734 & 0.024 \\
High emergency pressure & Supermajority required to invalidate laws & 0.551 & 0.577 & 0.443 & 0.643 & 0.738 & 0.023 \\
High emergency pressure & Strict recusal with substitute justices & 0.569 & 0.589 & 0.431 & 0.644 & 0.737 & 0.026 \\
High emergency pressure & Reasoned emergency orders with merits follow-up & 0.604 & 0.632 & 0.446 & 0.644 & 0.723 & 0.033 \\
High emergency pressure & Panel review with en banc safety valve & 0.559 & 0.587 & 0.430 & 0.644 & 0.736 & 0.024 \\
High emergency pressure & Dual supreme courts with cross-checking invalidations & 0.596 & 0.624 & 0.596 & 0.640 & 0.738 & 0.037 \\
High emergency pressure & Pre-enactment constitutional council & 0.565 & 0.584 & 0.661 & 0.643 & 0.743 & 0.027 \\
High emergency pressure & Judicial review with delayed legislative override & 0.540 & 0.580 & 0.432 & 0.644 & 0.740 & 0.024 \\
High emergency pressure & Retention-election accountability court & 0.562 & 0.598 & 0.408 & 0.644 & 0.741 & 0.029 \\
High emergency pressure & Hybrid court balancing independence and accountability & 0.612 & 0.640 & 0.432 & 0.641 & 0.736 & 0.039 \\
High emergency pressure & Weak-form review with legislative reply & 0.611 & 0.618 & 0.229 & 0.678 & 0.777 & 0.030 \\
High emergency pressure & Suspended declarations of invalidity & 0.648 & 0.660 & 0.247 & 0.623 & 0.841 & 0.035 \\
High emergency pressure & Strong-form review with explicit override clause & 0.552 & 0.590 & 0.373 & 0.632 & 0.880 & 0.025 \\
High emergency pressure & Pre-enactment review before laws take effect & 0.583 & 0.595 & 0.549 & 0.580 & 0.850 & 0.029 \\
High emergency pressure & Abstract review tribunal & 0.610 & 0.636 & 0.356 & 0.564 & 0.911 & 0.033 \\
High emergency pressure & Ombudsman-triggered constitutional review & 0.558 & 0.576 & 0.521 & 0.604 & 0.792 & 0.024 \\
High emergency pressure & Constitutional public defender access model & 0.599 & 0.631 & 0.514 & 0.632 & 0.750 & 0.030 \\
High emergency pressure & Rights-impact statements before review & 0.613 & 0.617 & 0.539 & 0.699 & 0.727 & 0.031 \\
High emergency pressure & Mandatory legislative response cycles & 0.628 & 0.655 & 0.265 & 0.615 & 0.920 & 0.047 \\
Low appointment polarization & Current-style federal supreme court & 0.539 & 0.586 & 0.528 & 0.646 & 0.682 & 0.031 \\
Low appointment polarization & Staggered 18-year nonrenewable terms & 0.549 & 0.593 & 0.519 & 0.646 & 0.681 & 0.033 \\
Low appointment polarization & Fifteen-justice supermajority commission court & 0.625 & 0.653 & 0.436 & 0.646 & 0.672 & 0.061 \\
Low appointment polarization & Supermajority required to invalidate laws & 0.623 & 0.654 & 0.445 & 0.645 & 0.675 & 0.064 \\
Low appointment polarization & Strict recusal with substitute justices & 0.637 & 0.661 & 0.434 & 0.646 & 0.675 & 0.077 \\
Low appointment polarization & Reasoned emergency orders with merits follow-up & 0.643 & 0.664 & 0.445 & 0.646 & 0.673 & 0.076 \\
Low appointment polarization & Panel review with en banc safety valve & 0.624 & 0.656 & 0.433 & 0.645 & 0.674 & 0.070 \\
Low appointment polarization & Dual supreme courts with cross-checking invalidations & 0.641 & 0.662 & 0.596 & 0.642 & 0.688 & 0.087 \\
Low appointment polarization & Pre-enactment constitutional council & 0.648 & 0.664 & 0.662 & 0.645 & 0.681 & 0.098 \\
Low appointment polarization & Judicial review with delayed legislative override & 0.609 & 0.652 & 0.434 & 0.646 & 0.678 & 0.063 \\
Low appointment polarization & Retention-election accountability court & 0.637 & 0.670 & 0.407 & 0.645 & 0.679 & 0.086 \\
Low appointment polarization & Hybrid court balancing independence and accountability & 0.636 & 0.675 & 0.432 & 0.643 & 0.687 & 0.091 \\
Low appointment polarization & Weak-form review with legislative reply & 0.705 & 0.718 & 0.179 & 0.680 & 0.714 & 0.281 \\
Low appointment polarization & Suspended declarations of invalidity & 0.689 & 0.692 & 0.211 & 0.626 & 0.791 & 0.075 \\
Low appointment polarization & Strong-form review with explicit override clause & 0.629 & 0.668 & 0.357 & 0.634 & 0.817 & 0.084 \\
Low appointment polarization & Pre-enactment review before laws take effect & 0.666 & 0.678 & 0.511 & 0.582 & 0.787 & 0.101 \\
Low appointment polarization & Abstract review tribunal & 0.655 & 0.676 & 0.310 & 0.567 & 0.859 & 0.081 \\
Low appointment polarization & Ombudsman-triggered constitutional review & 0.621 & 0.644 & 0.525 & 0.606 & 0.730 & 0.065 \\
Low appointment polarization & Constitutional public defender access model & 0.631 & 0.656 & 0.513 & 0.634 & 0.700 & 0.063 \\
Low appointment polarization & Rights-impact statements before review & 0.701 & 0.703 & 0.501 & 0.701 & 0.664 & 0.143 \\
Low appointment polarization & Mandatory legislative response cycles & 0.674 & 0.692 & 0.228 & 0.618 & 0.871 & 0.123 \\
High appointment polarization & Current-style federal supreme court & 0.502 & 0.550 & 0.532 & 0.642 & 0.781 & 0.023 \\
High appointment polarization & Staggered 18-year nonrenewable terms & 0.512 & 0.557 & 0.523 & 0.643 & 0.780 & 0.025 \\
High appointment polarization & Fifteen-justice supermajority commission court & 0.595 & 0.628 & 0.438 & 0.643 & 0.769 & 0.046 \\
High appointment polarization & Supermajority required to invalidate laws & 0.580 & 0.620 & 0.448 & 0.642 & 0.773 & 0.042 \\
High appointment polarization & Strict recusal with substitute justices & 0.603 & 0.633 & 0.436 & 0.642 & 0.772 & 0.051 \\
High appointment polarization & Reasoned emergency orders with merits follow-up & 0.601 & 0.630 & 0.450 & 0.642 & 0.771 & 0.046 \\
High appointment polarization & Panel review with en banc safety valve & 0.592 & 0.630 & 0.435 & 0.642 & 0.771 & 0.046 \\
High appointment polarization & Dual supreme courts with cross-checking invalidations & 0.610 & 0.636 & 0.597 & 0.639 & 0.785 & 0.057 \\
High appointment polarization & Pre-enactment constitutional council & 0.613 & 0.636 & 0.663 & 0.641 & 0.778 & 0.061 \\
High appointment polarization & Judicial review with delayed legislative override & 0.569 & 0.619 & 0.437 & 0.642 & 0.776 & 0.039 \\
High appointment polarization & Retention-election accountability court & 0.594 & 0.635 & 0.406 & 0.642 & 0.777 & 0.053 \\
High appointment polarization & Hybrid court balancing independence and accountability & 0.620 & 0.648 & 0.434 & 0.640 & 0.784 & 0.062 \\
High appointment polarization & Weak-form review with legislative reply & 0.689 & 0.699 & 0.169 & 0.676 & 0.810 & 0.134 \\
High appointment polarization & Suspended declarations of invalidity & 0.656 & 0.668 & 0.201 & 0.622 & 0.889 & 0.051 \\
High appointment polarization & Strong-form review with explicit override clause & 0.588 & 0.635 & 0.350 & 0.630 & 0.914 & 0.052 \\
High appointment polarization & Pre-enactment review before laws take effect & 0.635 & 0.654 & 0.512 & 0.579 & 0.884 & 0.069 \\
High appointment polarization & Abstract review tribunal & 0.623 & 0.651 & 0.314 & 0.563 & 0.956 & 0.056 \\
High appointment polarization & Ombudsman-triggered constitutional review & 0.591 & 0.620 & 0.527 & 0.602 & 0.827 & 0.046 \\
High appointment polarization & Constitutional public defender access model & 0.601 & 0.633 & 0.516 & 0.630 & 0.797 & 0.045 \\
High appointment polarization & Rights-impact statements before review & 0.663 & 0.673 & 0.503 & 0.697 & 0.761 & 0.073 \\
High appointment polarization & Mandatory legislative response cycles & 0.639 & 0.664 & 0.218 & 0.614 & 0.966 & 0.073 \\
Low rights threat & Current-style federal supreme court & 0.553 & 0.602 & 0.544 & 0.644 & 0.723 & 0.039 \\
Low rights threat & Staggered 18-year nonrenewable terms & 0.565 & 0.610 & 0.535 & 0.645 & 0.722 & 0.041 \\
Low rights threat & Fifteen-justice supermajority commission court & 0.659 & 0.682 & 0.452 & 0.645 & 0.712 & 0.101 \\
Low rights threat & Supermajority required to invalidate laws & 0.650 & 0.678 & 0.461 & 0.644 & 0.716 & 0.101 \\
Low rights threat & Strict recusal with substitute justices & 0.666 & 0.688 & 0.450 & 0.644 & 0.715 & 0.120 \\
Low rights threat & Reasoned emergency orders with merits follow-up & 0.673 & 0.690 & 0.461 & 0.644 & 0.714 & 0.140 \\
Low rights threat & Panel review with en banc safety valve & 0.657 & 0.684 & 0.449 & 0.644 & 0.715 & 0.115 \\
Low rights threat & Dual supreme courts with cross-checking invalidations & 0.670 & 0.688 & 0.612 & 0.641 & 0.728 & 0.186 \\
Low rights threat & Pre-enactment constitutional council & 0.668 & 0.684 & 0.677 & 0.643 & 0.722 & 0.182 \\
Low rights threat & Judicial review with delayed legislative override & 0.637 & 0.677 & 0.450 & 0.644 & 0.719 & 0.092 \\
Low rights threat & Retention-election accountability court & 0.660 & 0.693 & 0.425 & 0.644 & 0.720 & 0.137 \\
Low rights threat & Hybrid court balancing independence and accountability & 0.683 & 0.701 & 0.448 & 0.642 & 0.727 & 0.220 \\
Low rights threat & Weak-form review with legislative reply & 0.702 & 0.726 & 0.221 & 0.678 & 0.756 & 0.338 \\
Low rights threat & Suspended declarations of invalidity & 0.713 & 0.713 & 0.255 & 0.624 & 0.830 & 0.145 \\
Low rights threat & Strong-form review with explicit override clause & 0.652 & 0.689 & 0.389 & 0.632 & 0.858 & 0.134 \\
Low rights threat & Pre-enactment review before laws take effect & 0.685 & 0.697 & 0.526 & 0.581 & 0.828 & 0.201 \\
Low rights threat & Abstract review tribunal & 0.686 & 0.702 & 0.326 & 0.565 & 0.899 & 0.145 \\
Low rights threat & Ombudsman-triggered constitutional review & 0.657 & 0.674 & 0.539 & 0.604 & 0.770 & 0.110 \\
Low rights threat & Constitutional public defender access model & 0.665 & 0.682 & 0.527 & 0.632 & 0.740 & 0.120 \\
Low rights threat & Rights-impact statements before review & 0.719 & 0.720 & 0.517 & 0.699 & 0.705 & 0.297 \\
Low rights threat & Mandatory legislative response cycles & 0.693 & 0.712 & 0.270 & 0.616 & 0.913 & 0.247 \\
High rights threat & Current-style federal supreme court & 0.499 & 0.561 & 0.494 & 0.644 & 0.738 & 0.020 \\
High rights threat & Staggered 18-year nonrenewable terms & 0.505 & 0.564 & 0.486 & 0.645 & 0.737 & 0.020 \\
High rights threat & Fifteen-justice supermajority commission court & 0.572 & 0.629 & 0.401 & 0.645 & 0.726 & 0.030 \\
High rights threat & Supermajority required to invalidate laws & 0.566 & 0.625 & 0.411 & 0.644 & 0.730 & 0.028 \\
High rights threat & Strict recusal with substitute justices & 0.577 & 0.631 & 0.400 & 0.644 & 0.729 & 0.030 \\
High rights threat & Reasoned emergency orders with merits follow-up & 0.584 & 0.640 & 0.409 & 0.644 & 0.728 & 0.030 \\
High rights threat & Panel review with en banc safety valve & 0.565 & 0.630 & 0.398 & 0.644 & 0.729 & 0.030 \\
High rights threat & Dual supreme courts with cross-checking invalidations & 0.574 & 0.629 & 0.563 & 0.641 & 0.742 & 0.032 \\
High rights threat & Pre-enactment constitutional council & 0.572 & 0.629 & 0.630 & 0.643 & 0.736 & 0.034 \\
High rights threat & Judicial review with delayed legislative override & 0.553 & 0.626 & 0.401 & 0.644 & 0.732 & 0.028 \\
High rights threat & Retention-election accountability court & 0.561 & 0.625 & 0.373 & 0.644 & 0.735 & 0.033 \\
High rights threat & Hybrid court balancing independence and accountability & 0.585 & 0.643 & 0.398 & 0.641 & 0.742 & 0.034 \\
High rights threat & Weak-form review with legislative reply & 0.690 & 0.700 & 0.085 & 0.678 & 0.766 & 0.078 \\
High rights threat & Suspended declarations of invalidity & 0.641 & 0.677 & 0.116 & 0.624 & 0.847 & 0.032 \\
High rights threat & Strong-form review with explicit override clause & 0.572 & 0.641 & 0.295 & 0.632 & 0.869 & 0.036 \\
High rights threat & Pre-enactment review before laws take effect & 0.596 & 0.644 & 0.481 & 0.580 & 0.843 & 0.036 \\
High rights threat & Abstract review tribunal & 0.598 & 0.653 & 0.278 & 0.565 & 0.914 & 0.033 \\
High rights threat & Ombudsman-triggered constitutional review & 0.565 & 0.621 & 0.494 & 0.603 & 0.786 & 0.029 \\
High rights threat & Constitutional public defender access model & 0.581 & 0.650 & 0.479 & 0.631 & 0.756 & 0.031 \\
High rights threat & Rights-impact statements before review & 0.622 & 0.663 & 0.471 & 0.699 & 0.719 & 0.038 \\
High rights threat & Mandatory legislative response cycles & 0.613 & 0.663 & 0.132 & 0.616 & 0.921 & 0.044 \\
High public trust & Current-style federal supreme court & 0.544 & 0.594 & 0.526 & 0.661 & 0.692 & 0.067 \\
High public trust & Staggered 18-year nonrenewable terms & 0.557 & 0.603 & 0.517 & 0.662 & 0.691 & 0.075 \\
High public trust & Fifteen-justice supermajority commission court & 0.654 & 0.682 & 0.434 & 0.662 & 0.680 & 0.178 \\
High public trust & Supermajority required to invalidate laws & 0.645 & 0.679 & 0.443 & 0.661 & 0.684 & 0.171 \\
High public trust & Strict recusal with substitute justices & 0.663 & 0.689 & 0.432 & 0.661 & 0.683 & 0.213 \\
High public trust & Reasoned emergency orders with merits follow-up & 0.666 & 0.689 & 0.443 & 0.661 & 0.682 & 0.196 \\
High public trust & Panel review with en banc safety valve & 0.651 & 0.684 & 0.431 & 0.661 & 0.683 & 0.179 \\
High public trust & Dual supreme courts with cross-checking invalidations & 0.664 & 0.689 & 0.594 & 0.658 & 0.696 & 0.290 \\
High public trust & Pre-enactment constitutional council & 0.662 & 0.687 & 0.660 & 0.660 & 0.690 & 0.315 \\
High public trust & Judicial review with delayed legislative override & 0.636 & 0.678 & 0.433 & 0.661 & 0.687 & 0.170 \\
High public trust & Retention-election accountability court & 0.657 & 0.692 & 0.406 & 0.661 & 0.688 & 0.251 \\
High public trust & Hybrid court balancing independence and accountability & 0.678 & 0.701 & 0.430 & 0.659 & 0.695 & 0.309 \\
High public trust & Weak-form review with legislative reply & 0.719 & 0.740 & 0.177 & 0.695 & 0.723 & 0.808 \\
High public trust & Suspended declarations of invalidity & 0.719 & 0.724 & 0.208 & 0.641 & 0.799 & 0.230 \\
High public trust & Strong-form review with explicit override clause & 0.656 & 0.695 & 0.348 & 0.649 & 0.825 & 0.270 \\
High public trust & Pre-enactment review before laws take effect & 0.686 & 0.702 & 0.508 & 0.598 & 0.796 & 0.374 \\
High public trust & Abstract review tribunal & 0.682 & 0.706 & 0.309 & 0.582 & 0.868 & 0.247 \\
High public trust & Ombudsman-triggered constitutional review & 0.646 & 0.672 & 0.523 & 0.621 & 0.738 & 0.166 \\
High public trust & Constitutional public defender access model & 0.658 & 0.685 & 0.511 & 0.649 & 0.709 & 0.189 \\
High public trust & Rights-impact statements before review & 0.714 & 0.723 & 0.499 & 0.716 & 0.673 & 0.451 \\
High public trust & Mandatory legislative response cycles & 0.695 & 0.719 & 0.225 & 0.633 & 0.879 & 0.496 \\
Low public trust & Current-style federal supreme court & 0.504 & 0.542 & 0.534 & 0.620 & 0.785 & 0.003 \\
Low public trust & Staggered 18-year nonrenewable terms & 0.512 & 0.547 & 0.525 & 0.620 & 0.784 & 0.003 \\
Low public trust & Fifteen-justice supermajority commission court & 0.581 & 0.606 & 0.442 & 0.620 & 0.774 & 0.006 \\
Low public trust & Supermajority required to invalidate laws & 0.575 & 0.603 & 0.450 & 0.619 & 0.777 & 0.005 \\
Low public trust & Strict recusal with substitute justices & 0.586 & 0.609 & 0.440 & 0.620 & 0.777 & 0.005 \\
Low public trust & Reasoned emergency orders with merits follow-up & 0.593 & 0.612 & 0.452 & 0.619 & 0.776 & 0.006 \\
Low public trust & Panel review with en banc safety valve & 0.577 & 0.605 & 0.439 & 0.619 & 0.777 & 0.005 \\
Low public trust & Dual supreme courts with cross-checking invalidations & 0.586 & 0.607 & 0.601 & 0.616 & 0.790 & 0.006 \\
Low public trust & Pre-enactment constitutional council & 0.587 & 0.606 & 0.667 & 0.619 & 0.783 & 0.007 \\
Low public trust & Judicial review with delayed legislative override & 0.562 & 0.601 & 0.439 & 0.620 & 0.781 & 0.005 \\
Low public trust & Retention-election accountability court & 0.575 & 0.609 & 0.410 & 0.619 & 0.782 & 0.006 \\
Low public trust & Hybrid court balancing independence and accountability & 0.599 & 0.620 & 0.438 & 0.617 & 0.789 & 0.006 \\
Low public trust & Weak-form review with legislative reply & 0.645 & 0.654 & 0.178 & 0.654 & 0.815 & 0.011 \\
Low public trust & Suspended declarations of invalidity & 0.631 & 0.635 & 0.207 & 0.599 & 0.894 & 0.005 \\
Low public trust & Strong-form review with explicit override clause & 0.572 & 0.611 & 0.371 & 0.608 & 0.920 & 0.006 \\
Low public trust & Pre-enactment review before laws take effect & 0.604 & 0.618 & 0.516 & 0.556 & 0.889 & 0.007 \\
Low public trust & Abstract review tribunal & 0.601 & 0.624 & 0.316 & 0.540 & 0.961 & 0.006 \\
Low public trust & Ombudsman-triggered constitutional review & 0.576 & 0.597 & 0.531 & 0.579 & 0.832 & 0.006 \\
Low public trust & Constitutional public defender access model & 0.587 & 0.611 & 0.520 & 0.607 & 0.803 & 0.005 \\
Low public trust & Rights-impact statements before review & 0.635 & 0.640 & 0.506 & 0.675 & 0.766 & 0.008 \\
Low public trust & Mandatory legislative response cycles & 0.612 & 0.633 & 0.224 & 0.591 & 0.972 & 0.007 \\
Low legislative conflict & Current-style federal supreme court & 0.540 & 0.585 & 0.527 & 0.646 & 0.680 & 0.031 \\
Low legislative conflict & Staggered 18-year nonrenewable terms & 0.549 & 0.591 & 0.518 & 0.647 & 0.679 & 0.031 \\
Low legislative conflict & Fifteen-justice supermajority commission court & 0.647 & 0.671 & 0.434 & 0.647 & 0.668 & 0.072 \\
Low legislative conflict & Supermajority required to invalidate laws & 0.636 & 0.665 & 0.443 & 0.646 & 0.672 & 0.067 \\
Low legislative conflict & Strict recusal with substitute justices & 0.653 & 0.675 & 0.431 & 0.646 & 0.671 & 0.081 \\
Low legislative conflict & Reasoned emergency orders with merits follow-up & 0.656 & 0.675 & 0.443 & 0.646 & 0.670 & 0.086 \\
Low legislative conflict & Panel review with en banc safety valve & 0.642 & 0.670 & 0.431 & 0.646 & 0.671 & 0.075 \\
Low legislative conflict & Dual supreme courts with cross-checking invalidations & 0.655 & 0.675 & 0.594 & 0.642 & 0.684 & 0.102 \\
Low legislative conflict & Pre-enactment constitutional council & 0.658 & 0.674 & 0.660 & 0.645 & 0.678 & 0.114 \\
Low legislative conflict & Judicial review with delayed legislative override & 0.625 & 0.663 & 0.433 & 0.646 & 0.675 & 0.062 \\
Low legislative conflict & Retention-election accountability court & 0.653 & 0.680 & 0.407 & 0.646 & 0.676 & 0.091 \\
Low legislative conflict & Hybrid court balancing independence and accountability & 0.667 & 0.687 & 0.430 & 0.644 & 0.683 & 0.100 \\
Low legislative conflict & Weak-form review with legislative reply & 0.713 & 0.722 & 0.177 & 0.680 & 0.711 & 0.288 \\
Low legislative conflict & Suspended declarations of invalidity & 0.711 & 0.711 & 0.209 & 0.626 & 0.787 & 0.089 \\
Low legislative conflict & Strong-form review with explicit override clause & 0.650 & 0.683 & 0.348 & 0.634 & 0.813 & 0.097 \\
Low legislative conflict & Pre-enactment review before laws take effect & 0.681 & 0.689 & 0.508 & 0.582 & 0.784 & 0.130 \\
Low legislative conflict & Abstract review tribunal & 0.673 & 0.692 & 0.309 & 0.567 & 0.856 & 0.094 \\
Low legislative conflict & Ombudsman-triggered constitutional review & 0.641 & 0.660 & 0.522 & 0.606 & 0.726 & 0.073 \\
Low legislative conflict & Constitutional public defender access model & 0.649 & 0.671 & 0.510 & 0.634 & 0.697 & 0.074 \\
Low legislative conflict & Rights-impact statements before review & 0.708 & 0.709 & 0.499 & 0.701 & 0.661 & 0.139 \\
Low legislative conflict & Mandatory legislative response cycles & 0.690 & 0.705 & 0.226 & 0.618 & 0.867 & 0.165 \\
High legislative conflict & Current-style federal supreme court & 0.501 & 0.544 & 0.535 & 0.628 & 0.816 & 0.009 \\
High legislative conflict & Staggered 18-year nonrenewable terms & 0.511 & 0.550 & 0.526 & 0.628 & 0.816 & 0.010 \\
High legislative conflict & Fifteen-justice supermajority commission court & 0.577 & 0.606 & 0.443 & 0.628 & 0.806 & 0.014 \\
High legislative conflict & Supermajority required to invalidate laws & 0.572 & 0.603 & 0.451 & 0.627 & 0.809 & 0.015 \\
High legislative conflict & Strict recusal with substitute justices & 0.582 & 0.610 & 0.441 & 0.628 & 0.809 & 0.015 \\
High legislative conflict & Reasoned emergency orders with merits follow-up & 0.589 & 0.613 & 0.453 & 0.627 & 0.808 & 0.016 \\
High legislative conflict & Panel review with en banc safety valve & 0.572 & 0.607 & 0.440 & 0.627 & 0.809 & 0.015 \\
High legislative conflict & Dual supreme courts with cross-checking invalidations & 0.583 & 0.607 & 0.602 & 0.624 & 0.822 & 0.018 \\
High legislative conflict & Pre-enactment constitutional council & 0.583 & 0.606 & 0.668 & 0.626 & 0.815 & 0.019 \\
High legislative conflict & Judicial review with delayed legislative override & 0.556 & 0.601 & 0.439 & 0.627 & 0.814 & 0.014 \\
High legislative conflict & Retention-election accountability court & 0.570 & 0.608 & 0.409 & 0.627 & 0.815 & 0.017 \\
High legislative conflict & Hybrid court balancing independence and accountability & 0.594 & 0.621 & 0.439 & 0.625 & 0.822 & 0.019 \\
High legislative conflict & Weak-form review with legislative reply & 0.638 & 0.651 & 0.176 & 0.662 & 0.847 & 0.029 \\
High legislative conflict & Suspended declarations of invalidity & 0.623 & 0.631 & 0.206 & 0.607 & 0.926 & 0.016 \\
High legislative conflict & Strong-form review with explicit override clause & 0.565 & 0.610 & 0.379 & 0.616 & 0.952 & 0.016 \\
High legislative conflict & Pre-enactment review before laws take effect & 0.600 & 0.619 & 0.516 & 0.564 & 0.921 & 0.020 \\
High legislative conflict & Abstract review tribunal & 0.597 & 0.624 & 0.317 & 0.548 & 0.989 & 0.017 \\
High legislative conflict & Ombudsman-triggered constitutional review & 0.573 & 0.597 & 0.532 & 0.587 & 0.864 & 0.015 \\
High legislative conflict & Constitutional public defender access model & 0.584 & 0.612 & 0.521 & 0.615 & 0.835 & 0.015 \\
High legislative conflict & Rights-impact statements before review & 0.631 & 0.639 & 0.508 & 0.683 & 0.798 & 0.021 \\
High legislative conflict & Mandatory legislative response cycles & 0.609 & 0.632 & 0.222 & 0.599 & 0.995 & 0.019 \\
\bottomrule
\end{longtable}
}


\section*{Legislative-Output Docket Stress}
The main manuscript treats the legislative-output bridge as a robustness exercise. Imported rows change docket stress fields such as rights threat, public support, legislative mandate, policy-domain mix, capture pressure, volatility, and emergency pressure; they do not change the constitutional-review rules. Figure~\ref{fig:paired-import} reports the directional display-score movement under these docket sources.

\begin{figure}[hbt!]
\centering
% Auto-generated by paper/scripts/generate_figures.py
\begingroup
\setlength{\unitlength}{1mm}
\begin{picture}(128,88)
\scriptsize
\put(24.0,14.0){\color{black!15}\line(1,0){82.0}}
\put(20.0,14.0){\makebox(0,0)[r]{0.50}}
\put(24.0,31.1){\color{black!15}\line(1,0){82.0}}
\put(20.0,31.1){\makebox(0,0)[r]{0.60}}
\put(24.0,48.1){\color{black!15}\line(1,0){82.0}}
\put(20.0,48.1){\makebox(0,0)[r]{0.70}}
\put(24.0,65.2){\color{black!15}\line(1,0){82.0}}
\put(20.0,65.2){\makebox(0,0)[r]{0.80}}
\put(24.0,14.0){\color{black!12}\line(0,1){58.0}}
\put(24.0,10.0){\makebox(0,0){SB}}
\put(44.5,14.0){\color{black!12}\line(0,1){58.0}}
\put(44.5,10.0){\makebox(0,0){ALL}}
\put(65.0,14.0){\color{black!12}\line(0,1){58.0}}
\put(65.0,10.0){\makebox(0,0){CAP}}
\put(85.5,14.0){\color{black!12}\line(0,1){58.0}}
\put(85.5,10.0){\makebox(0,0){VOL}}
\put(106.0,14.0){\color{black!12}\line(0,1){58.0}}
\put(106.0,10.0){\makebox(0,0){MAND}}
\put(24.0,14.0){\line(1,0){82.0}}
\put(24.0,14.0){\line(0,1){58.0}}
\linethickness{0.35mm}
{\color{black}\qbezier[22](24.0,13.3)(34.2,24.7)(44.5,36.2)}
{\color{black}\qbezier[22](44.5,36.2)(54.8,33.1)(65.0,30.0)}
{\color{black}\qbezier[22](65.0,30.0)(75.2,26.6)(85.5,23.2)}
{\color{black}\qbezier[22](85.5,23.2)(95.8,24.4)(106.0,25.6)}
\linethickness{0.25mm}
\put(24.0,13.3){\makebox(0,0){\color{black}\rule{1.5mm}{1.5mm}}}
\put(44.5,36.2){\makebox(0,0){\color{black}\rule{1.5mm}{1.5mm}}}
\put(65.0,30.0){\makebox(0,0){\color{black}\rule{1.5mm}{1.5mm}}}
\put(85.5,23.2){\makebox(0,0){\color{black}\rule{1.5mm}{1.5mm}}}
\put(106.0,25.6){\makebox(0,0){\color{black}\rule{1.5mm}{1.5mm}}}
\linethickness{0.35mm}
{\color{black!70}\qbezier[22](24.0,32.3)(34.2,40.7)(44.5,49.1)}
{\color{black!70}\qbezier[22](44.5,49.1)(54.8,48.9)(65.0,48.6)}
{\color{black!70}\qbezier[22](65.0,48.6)(75.2,43.7)(85.5,38.7)}
{\color{black!70}\qbezier[22](85.5,38.7)(95.8,39.7)(106.0,40.6)}
\linethickness{0.25mm}
\put(24.0,32.3){\makebox(0,0){\color{black!70}\rule{1.5mm}{1.5mm}}}
\put(44.5,49.1){\makebox(0,0){\color{black!70}\rule{1.5mm}{1.5mm}}}
\put(65.0,48.6){\makebox(0,0){\color{black!70}\rule{1.5mm}{1.5mm}}}
\put(85.5,38.7){\makebox(0,0){\color{black!70}\rule{1.5mm}{1.5mm}}}
\put(106.0,40.6){\makebox(0,0){\color{black!70}\rule{1.5mm}{1.5mm}}}
\linethickness{0.35mm}
{\color{black!45}\qbezier[22](24.0,31.6)(34.2,39.3)(44.5,47.1)}
{\color{black!45}\qbezier[22](44.5,47.1)(54.8,42.3)(65.0,37.5)}
{\color{black!45}\qbezier[22](65.0,37.5)(75.2,38.2)(85.5,38.9)}
{\color{black!45}\qbezier[22](85.5,38.9)(95.8,36.3)(106.0,33.8)}
\linethickness{0.25mm}
\put(24.0,31.6){\makebox(0,0){\color{black!45}\rule{1.5mm}{1.5mm}}}
\put(44.5,47.1){\makebox(0,0){\color{black!45}\rule{1.5mm}{1.5mm}}}
\put(65.0,37.5){\makebox(0,0){\color{black!45}\rule{1.5mm}{1.5mm}}}
\put(85.5,38.9){\makebox(0,0){\color{black!45}\rule{1.5mm}{1.5mm}}}
\put(106.0,33.8){\makebox(0,0){\color{black!45}\rule{1.5mm}{1.5mm}}}
\linethickness{0.35mm}
{\color{black!30}\qbezier[22](24.0,29.7)(34.2,38.4)(44.5,47.1)}
{\color{black!30}\qbezier[22](44.5,47.1)(54.8,46.9)(65.0,46.8)}
{\color{black!30}\qbezier[22](65.0,46.8)(75.2,43.9)(85.5,41.1)}
{\color{black!30}\qbezier[22](85.5,41.1)(95.8,42.4)(106.0,43.7)}
\linethickness{0.25mm}
\put(24.0,29.7){\makebox(0,0){\color{black!30}\rule{1.5mm}{1.5mm}}}
\put(44.5,47.1){\makebox(0,0){\color{black!30}\rule{1.5mm}{1.5mm}}}
\put(65.0,46.8){\makebox(0,0){\color{black!30}\rule{1.5mm}{1.5mm}}}
\put(85.5,41.1){\makebox(0,0){\color{black!30}\rule{1.5mm}{1.5mm}}}
\put(106.0,43.7){\makebox(0,0){\color{black!30}\rule{1.5mm}{1.5mm}}}
\linethickness{0.35mm}
{\color{black!55}\qbezier[22](24.0,33.8)(34.2,41.8)(44.5,49.8)}
{\color{black!55}\qbezier[22](44.5,49.8)(54.8,48.7)(65.0,47.6)}
{\color{black!55}\qbezier[22](65.0,47.6)(75.2,44.5)(85.5,41.3)}
{\color{black!55}\qbezier[22](85.5,41.3)(95.8,40.8)(106.0,40.3)}
\linethickness{0.25mm}
\put(24.0,33.8){\makebox(0,0){\color{black!55}\rule{1.5mm}{1.5mm}}}
\put(44.5,49.8){\makebox(0,0){\color{black!55}\rule{1.5mm}{1.5mm}}}
\put(65.0,47.6){\makebox(0,0){\color{black!55}\rule{1.5mm}{1.5mm}}}
\put(85.5,41.3){\makebox(0,0){\color{black!55}\rule{1.5mm}{1.5mm}}}
\put(106.0,40.3){\makebox(0,0){\color{black!55}\rule{1.5mm}{1.5mm}}}
\linethickness{0.35mm}
{\color{black!40}\qbezier[22](24.0,39.9)(34.2,45.8)(44.5,51.7)}
{\color{black!40}\qbezier[22](44.5,51.7)(54.8,50.1)(65.0,48.5)}
{\color{black!40}\qbezier[22](65.0,48.5)(75.2,48.3)(85.5,48.1)}
{\color{black!40}\qbezier[22](85.5,48.1)(95.8,49.3)(106.0,50.5)}
\linethickness{0.25mm}
\put(24.0,39.9){\makebox(0,0){\color{black!40}\rule{1.5mm}{1.5mm}}}
\put(44.5,51.7){\makebox(0,0){\color{black!40}\rule{1.5mm}{1.5mm}}}
\put(65.0,48.5){\makebox(0,0){\color{black!40}\rule{1.5mm}{1.5mm}}}
\put(85.5,48.1){\makebox(0,0){\color{black!40}\rule{1.5mm}{1.5mm}}}
\put(106.0,50.5){\makebox(0,0){\color{black!40}\rule{1.5mm}{1.5mm}}}
\linethickness{0.35mm}
{\color{black!60}\qbezier[22](24.0,36.0)(34.2,42.7)(44.5,49.5)}
{\color{black!60}\qbezier[22](44.5,49.5)(54.8,49.3)(65.0,49.1)}
{\color{black!60}\qbezier[22](65.0,49.1)(75.2,46.1)(85.5,43.0)}
{\color{black!60}\qbezier[22](85.5,43.0)(95.8,40.1)(106.0,37.2)}
\linethickness{0.25mm}
\put(24.0,36.0){\makebox(0,0){\color{black!60}\rule{1.5mm}{1.5mm}}}
\put(44.5,49.5){\makebox(0,0){\color{black!60}\rule{1.5mm}{1.5mm}}}
\put(65.0,49.1){\makebox(0,0){\color{black!60}\rule{1.5mm}{1.5mm}}}
\put(85.5,43.0){\makebox(0,0){\color{black!60}\rule{1.5mm}{1.5mm}}}
\put(106.0,37.2){\makebox(0,0){\color{black!60}\rule{1.5mm}{1.5mm}}}
{\color{black!55}\qbezier[10](107.1,25.6)(108.0,25.1)(109.0,24.6)}
\put(110.0,24.6){\makebox(0,0)[l]{\color{black}CUR*}}
{\color{black!55}\qbezier[10](107.1,40.6)(108.0,45.9)(109.0,51.2)}
\put(110.0,51.2){\makebox(0,0)[l]{\color{black}EMR}}
{\color{black!55}\qbezier[10](107.1,33.8)(108.0,37.0)(109.0,40.2)}
\put(110.0,40.2){\makebox(0,0)[l]{\color{black}DUAL}}
{\color{black!55}\qbezier[10](107.1,43.7)(108.0,50.2)(109.0,56.7)}
\put(110.0,56.7){\makebox(0,0)[l]{\color{black}COUNC}}
{\color{black!55}\qbezier[10](107.1,40.3)(108.0,43.0)(109.0,45.7)}
\put(110.0,45.7){\makebox(0,0)[l]{\color{black}HYB}}
{\color{black!55}\qbezier[10](107.1,50.5)(108.0,56.4)(109.0,62.2)}
\put(110.0,62.2){\makebox(0,0)[l]{\color{black}WFR}}
{\color{black!55}\qbezier[10](107.1,37.2)(108.0,36.0)(109.0,34.7)}
\put(110.0,34.7){\makebox(0,0)[l]{\color{black}MLR}}
\put(65,4){\makebox(0,0){docket source}}
\put(12,75){\makebox(0,0){score}}
\put(65,80){\makebox(0,0){SB generated; ALL all imports; CAP high capture; VOL high volatility; MAND low mandate}}
\end{picture}
\endgroup

\caption{Directional display scores under generated and legislative-output docket sources. SB means generated baseline, ALL all imported rows, CAP high capture, VOL high volatility, and MAND low mandate. CUR* marks the current-style scenario.}
\label{fig:paired-import}
\end{figure}

\section*{Calibration Source Matrix}
The calibration matrix distinguishes source-backed checks from contextual stress-test ranges. Only rows marked for validation in the table are counted in the manuscript's source-range comparisons. Rows with synthesis labels, missing denominators, public-trust proxies, or normalized cost construction are retained for transparency but excluded from fit counts.
% Auto-generated by paper/scripts/generate_supplement.py
{\scriptsize
\setlength{\tabcolsep}{2pt}
\begin{longtable}{@{}p{0.17\linewidth}p{0.24\linewidth}p{0.20\linewidth}p{0.05\linewidth}p{0.08\linewidth}p{0.14\linewidth}@{}}
\caption{Calibration source-observation matrix.}\label{tab:supp-calibration-sources}\\
\toprule
Profile & Target & Method & $n$ & Rel. & Source status \\
\midrule
\endfirsthead
\toprule
Profile & Target & Method & $n$ & Rel. & Source status \\
\midrule
\endhead
canada-scc-recent & Leave application grant rate & leave applications and grants & 524 & high & URL+n \\
canada-scc-recent & Charter invalidation proxy & Charter outcome synthesis & -- & medium & synthesis \\
canada-scc-recent & Public trust proxy & public-trust survey proxy & -- & medium & synthesis \\
canada-scc-recent & Normalized direct court cost & benchmark cost profile & 9 & medium & counted \\
france-conseil-qpc & QPC invalidation rate & QPC decision outcome summary & 67 & medium & URL+n \\
france-conseil-qpc & Public trust proxy & public-trust survey proxy & -- & low & synthesis \\
france-conseil-qpc & Normalized direct court cost & benchmark cost profile & 9 & medium & counted \\
france-conseil-qpc & Normalized upstream screening cost & benchmark cost profile & -- & medium & synthesis \\
germany-bverfg-2024 & Constitutional complaint success and admission proxy & annual caseload and complaint-grant proxy & 4640 & medium & URL+n \\
germany-bverfg-2024 & Public trust proxy & public-trust survey proxy & -- & medium & synthesis \\
germany-bverfg-2024 & Normalized direct court cost & benchmark cost profile & 16 & medium & counted \\
germany-bverfg-2024 & Normalized capacity strain & benchmark cost profile & -- & medium & synthesis \\
cost-us-supreme-court & Normalized direct court cost & budget and staffing benchmark & 9 & medium & URL+n \\
cost-us-supreme-court & Normalized capacity strain & budget and filings benchmark & -- & medium & URL \\
cost-uk-supreme-court & Normalized direct court cost & budget and staffing benchmark & 12 & medium & URL+n \\
cost-france-conseil & Normalized direct court cost & budget and decision benchmark & 9 & medium & URL+n \\
south-africa-constcourt-recent & Merits invalidation proxy & merits outcome synthesis & -- & medium & URL \\
south-africa-constcourt-recent & Petition-to-judgment throughput proxy & annual petitions and judgments & 355 & medium & URL+n \\
south-africa-constcourt-recent & Public trust proxy & public-trust survey proxy & -- & medium & synthesis \\
south-africa-constcourt-recent & Normalized capacity strain & benchmark cost profile & -- & medium & synthesis \\
scdb-postwar-merits-1946-2024 & Speech docket share & SCDB issue-area proxy & -- & high & URL \\
scdb-postwar-merits-1946-2024 & Civil-rights and privacy docket share & SCDB issue-area proxy & -- & high & URL \\
scdb-postwar-merits-1946-2024 & Criminal procedure docket share & SCDB issue-area proxy & -- & high & URL \\
scdb-postwar-merits-1946-2024 & Federalism docket share & SCDB issue-area proxy & -- & high & URL \\
scdb-postwar-merits-1946-2024 & Election-law docket share & SCDB issue-code proxy & -- & medium & URL \\
scdb-postwar-merits-1946-2024 & Emergency-powers merits share & SCDB issue-code proxy & -- & medium & URL \\
scdb-postwar-merits-1946-2024 & Administrative and economic regulation share & SCDB issue-area plus agency-review proxy & -- & medium & URL \\
scdb-modern-merits-2000-2024 & Speech docket share & SCDB issue-area proxy & -- & high & URL \\
scdb-modern-merits-2000-2024 & Civil-rights and privacy docket share & SCDB issue-area proxy & -- & high & URL \\
scdb-modern-merits-2000-2024 & Criminal procedure docket share & SCDB issue-area proxy & -- & high & URL \\
scdb-modern-merits-2000-2024 & Federalism docket share & SCDB issue-area proxy & -- & high & URL \\
scdb-modern-merits-2000-2024 & Election-law docket share & SCDB issue-code proxy & -- & medium & URL \\
scdb-modern-merits-2000-2024 & Emergency-powers merits share & SCDB issue-code proxy & -- & medium & URL \\
scdb-modern-merits-2000-2024 & Administrative and economic regulation share & SCDB issue-area plus agency-review proxy & -- & medium & URL \\
scotus-emergency-2024-2025 & Substantive emergency application relief rate & SCOTUSblog emergency application universe summary & 43 & medium & URL+n \\
scotus-emergency-2024-2025 & Written explanation share & SCOTUSblog emergency application universe summary & 43 & medium & URL+n \\
scotus-emergency-2024-2025 & Public disagreement share & SCOTUSblog emergency application universe summary & 43 & medium & URL+n \\
scotus-emergency-2024-2025 & Certiorari or emergency screening acceptance proxy & OT2024 docketed matters and cert grants & 3854 & medium & URL+n \\
gallup-court-confidence-2024 & Public court trust and approval & Gallup confidence and approval bracket & -- & medium & URL \\
\bottomrule
\end{longtable}
\noindent\footnotesize Rows marked synthesis are included for transparent stress testing, not as source-backed validation claims. The CSV source retains notes, source names, URLs when available, and construction notes for every target.
}


\section*{Statistical Units and Source Evidence}
The object-level extension distinguishes statutes, regulations, and government conduct within a case. Its object counts and conditional relief allocation are explicitly synthetic. Only nullified statutes among merits statute dispositions enter the Canadian nullification check; weak-form declarations remain in the denominator but not the numerator, and emergency-only orders are excluded. The eight additional German and South African observations retain their source units without being promoted to unrelated model outputs.
% Auto-generated by paper/scripts/generate_supplement.py
{\footnotesize
\setlength{\tabcolsep}{3pt}
\begin{longtable}{@{}L{0.18\linewidth}L{0.36\linewidth}rrL{0.19\linewidth}@{}}
\caption{Denominator-backed measurement audit.}\label{tab:supp-measurement-audit}\\
\toprule Court / period & Source measure & Counts & Rate & Model use \\
\midrule\endfirsthead
\toprule Court / period & Source measure & Counts & Rate & Model use \\
\midrule\endhead
Canada 1982--89 & Statute nullification rate among federal and provincial statute dispositions & 19/50 & 38.00\% & Narrow range check \\
Germany 2024 & New complaints with legal representation & 1500/4436 & 33.81\% & Source context only \\
Germany 2024 & New complaints accompanied by interim applications & 1003/4436 & 22.61\% & Missing model measure \\
Germany 2024 & Successful constitutional complaint proceedings & 39/4595 & 0.85\% & Missing model measure \\
Germany 2024 & Direct challenges to laws and ordinances & 110/4436 & 2.48\% & Source context only \\
South Africa 2024 & Dismissed petitions with observable reasons & 40/57 & 70.18\% & Source context only \\
South Africa 2024 & New petitions dismissed by observation cutoff & 57/217 & 26.27\% & Source context only \\
South Africa 2024 & Petitions seeking direct access & 19/217 & 8.76\% & Source context only \\
South Africa 2024 & Petitions seeking leave to appeal & 189/217 & 87.10\% & Source context only \\
\bottomrule\end{longtable}
\noindent Canada counts statute dispositions, not decisions. German rows describe decided complaints, new filings, representation, interim applications, and challenged acts in separate universes. South African rows cover petitions filed from 1 January to 26 July 2024; the route subtotals leave one petition unaccounted for, dismissals are a censored snapshot, and missing orders are source missingness. These eight German and South African observations add verified evidence but no validation-counted rows. Full source URLs, locators, counts, and limitations are preserved in the measurement-audit CSV and profile cards.
}


\section*{Source-Range Miss Roadmap}
The main manuscript reports source-range misses when validation-counted checks fall outside documented ranges. When no such misses remain, the generated roadmap records that state while the empirical research roadmap continues to identify coverage-expansion tasks. A verified source measure remains excluded when its statistical unit is not exposed by the simulator; the platform records that condition as a model-metric gap rather than mapping the source to a broader output.
% Auto-generated by paper/scripts/generate_tables.py
\begin{longtable}{@{}L{0.20\linewidth}L{0.16\linewidth}L{0.16\linewidth}L{0.19\linewidth}L{0.10\linewidth}@{}}
\caption{Full source-range miss roadmap.}
\label{tab:supp-validation-misses}\\
\toprule
Target & Model vs. source & Category & Next action & Gap \\
\midrule
\endfirsthead
\toprule
Target & Model vs. source & Category & Next action & Gap \\
\midrule
\endhead
Court of Justice of the European Union: Full compliance with rule-of-law rulings & 0.642 vs. 0.534--0.632 & source-range check & Inspect whether the source denominator, simulator metric, and scenario preset are directly comparable. & 0.010 \\
French Constitutional Council: QPC nonconformity rate & 0.299 vs. 0.305--0.324 & merits-outcome mechanism & Retune doctrine severity, referral screening, and invalidation-threshold parameters only after preserving cross-profile comparability. & 0.006 \\
\bottomrule
\end{longtable}


\section*{Empirical Research Roadmap}
The empirical data-development files identify the exact data gaps that must be filled before the new response, access, transplant, and implementation measures can be promoted from stress-test assumptions to source-specific calibration targets. Candidate rows preserve useful numerators and denominators from source-discovery work, the generated source-acquisition packet turns roadmap-only gaps into collection briefs, and the source-promotion packet labels blockers in stress-only source rows; only rows marked as promoted and mirrored in the calibration source matrix enter source-range checks.
% Auto-generated by paper/scripts/generate_supplement.py
{\scriptsize
\setlength{\tabcolsep}{3pt}
\begin{longtable}{@{}L{0.24\linewidth}rL{0.48\linewidth}@{}}
\caption{Empirical research roadmap files.}\label{tab:supp-research-roadmap}\\
\toprule
Roadmap file & Rows & Purpose \\
\midrule
\endfirsthead
\toprule
Roadmap file & Rows & Purpose \\
\midrule
\endhead
empirical-target-roadmap.csv & 16 & Court calibration packs \\
legislative-response-evidence.csv & 10 & Legislative response \\
transplant-feasibility-factors.csv & 10 & Transplant factors \\
compliance-enforcement-targets.csv & 9 & Compliance channels \\
case-selection-targets.csv & 9 & Case selection \\
\bottomrule
\end{longtable}
\noindent\footnotesize Roadmap rows are source-gathering tasks, not calibration targets. A row should move into the calibration source matrix only after URLs, denominators, period coverage, coding rules, and a direct simulator analogue are documented.
}


\section*{Artifact Inventory}
The artifact inventory records generated report outputs available in the local analysis directory. The anonymous review bundle excludes raw case-level and interval outputs by default to keep the upload practical; those outputs are reproducible from the listed commands and should be included in the final public replication deposit.
% Auto-generated by paper/scripts/generate_supplement.py
\begin{longtable}{@{}L{0.66\linewidth}r@{}}
\caption{Generated report artifact inventory.}\label{tab:supp-artifacts}\\
\toprule
Artifact & Size (KB) \\
\midrule
\endfirsthead
\toprule
Artifact & Size (KB) \\
\midrule
\endhead
reports/constitutional-review-campaign-v0-calibration-intervals.csv & 12.3 \\
reports/constitutional-review-campaign-v0-calibration.csv & 17.5 \\
reports/constitutional-review-campaign-v0-composition-intervals.csv & 874.2 \\
reports/constitutional-review-campaign-v0-composition.csv & 116.6 \\
reports/constitutional-review-campaign-v0-doctrine-intervals.csv & 6485.5 \\
reports/constitutional-review-campaign-v0-doctrines.csv & 394.4 \\
reports/constitutional-review-campaign-v0-intervals.csv & 1170.6 \\
reports/constitutional-review-campaign-v0-period-intervals.csv & 3524.3 \\
reports/constitutional-review-campaign-v0-periods.csv & 220.5 \\
reports/constitutional-review-campaign-v0-pipeline-intervals.csv & 11583.2 \\
reports/constitutional-review-campaign-v0-pipelines.csv & 687.8 \\
reports/constitutional-review-campaign-v0-policy-domain-intervals.csv & 7568.8 \\
reports/constitutional-review-campaign-v0-policy-domains.csv & 455.1 \\
reports/constitutional-review-campaign-v0.csv & 73.8 \\
reports/constitutional-review-paired-import-v1-calibration-intervals.csv & 12.3 \\
reports/constitutional-review-paired-import-v1-calibration.csv & 17.5 \\
reports/constitutional-review-paired-import-v1-composition-intervals.csv & 913.6 \\
reports/constitutional-review-paired-import-v1-composition.csv & 121.5 \\
reports/constitutional-review-paired-import-v1-doctrine-intervals.csv & 6750.8 \\
reports/constitutional-review-paired-import-v1-doctrines.csv & 402.0 \\
reports/constitutional-review-paired-import-v1-intervals.csv & 1223.8 \\
reports/constitutional-review-paired-import-v1-period-intervals.csv & 3681.9 \\
reports/constitutional-review-paired-import-v1-periods.csv & 225.3 \\
reports/constitutional-review-paired-import-v1-pipeline-intervals.csv & 10437.4 \\
reports/constitutional-review-paired-import-v1-pipelines.csv & 608.4 \\
reports/constitutional-review-paired-import-v1-policy-domain-intervals.csv & 6074.5 \\
reports/constitutional-review-paired-import-v1-policy-domains.csv & 359.1 \\
reports/constitutional-review-paired-import-v1.csv & 74.6 \\
reports/constitutional-review-sensitivity-v1-calibration-intervals.csv & 12.3 \\
reports/constitutional-review-sensitivity-v1-calibration.csv & 17.5 \\
reports/constitutional-review-sensitivity-v1-composition-intervals.csv & 1766.7 \\
reports/constitutional-review-sensitivity-v1-composition.csv & 236.6 \\
reports/constitutional-review-sensitivity-v1-doctrine-intervals.csv & 13173.1 \\
reports/constitutional-review-sensitivity-v1-doctrines.csv & 829.8 \\
reports/constitutional-review-sensitivity-v1-intervals.csv & 2369.3 \\
reports/constitutional-review-sensitivity-v1-period-intervals.csv & 7155.6 \\
reports/constitutional-review-sensitivity-v1-periods.csv & 464.4 \\
reports/constitutional-review-sensitivity-v1-pipeline-intervals.csv & 23661.6 \\
reports/constitutional-review-sensitivity-v1-pipelines.csv & 1453.1 \\
reports/constitutional-review-sensitivity-v1-policy-domain-intervals.csv & 15423.8 \\
reports/constitutional-review-sensitivity-v1-policy-domains.csv & 959.6 \\
reports/constitutional-review-sensitivity-v1.csv & 155.4 \\
reports/constitutional-review-validation-v1-calibration-intervals.csv & 12.2 \\
reports/constitutional-review-validation-v1-calibration.csv & 17.4 \\
reports/constitutional-review-validation-v1-composition-intervals.csv & 46.7 \\
reports/constitutional-review-validation-v1-composition.csv & 6.4 \\
reports/constitutional-review-validation-v1-doctrine-intervals.csv & 344.0 \\
reports/constitutional-review-validation-v1-doctrines.csv & 21.0 \\
reports/constitutional-review-validation-v1-intervals.csv & 62.7 \\
reports/constitutional-review-validation-v1-period-intervals.csv & 188.0 \\
reports/constitutional-review-validation-v1-periods.csv & 12.2 \\
reports/constitutional-review-validation-v1-pipeline-intervals.csv & 611.9 \\
reports/constitutional-review-validation-v1-pipelines.csv & 35.8 \\
reports/constitutional-review-validation-v1-policy-domain-intervals.csv & 400.6 \\
reports/constitutional-review-validation-v1-policy-domains.csv & 24.1 \\
reports/constitutional-review-validation-v1.csv & 5.3 \\
reports/constitutional-review-campaign-v0-cases.csv.gz & 127460.5 \\
reports/constitutional-review-paired-import-v1-cases.csv.gz & 129006.4 \\
reports/constitutional-review-sensitivity-v1-cases.csv.gz & 186302.1 \\
reports/constitutional-review-validation-v1-cases.csv.gz & 6149.7 \\
\bottomrule
\end{longtable}


\end{document}
